% !TeX root = ../../Book1.tex
% 纲要标题：## 第10章　环、理想与中国剩余定理
% 纲要位置：计划纲要.md，第 386 行。

\chapter{环、理想与中国剩余定理}
\label{b1:ch:10}
\BookChapterNumber{b1:ch:10}

\begin{chapterabstract}
环把加法、乘法及分配律置于同一结构中；理想使这两种运算能够一起下降到商环。本章研究同态与理想、素理想和极大理想，以及理想间的运算，最后以中国剩余定理说明互素条件下的商环分解。
\end{chapterabstract}

\begin{learninggoals}
\begin{itemize}
\item 明确含幺环、保单位同态与子环的约定，区别单位、零因子和幂零元。
\item 判断左右理想与双边理想，构造商环并准确计算同态的核和像。
\item 用商环判断素理想与极大理想，比较理想的和、积、交、根与理想商。
\item 显式求解同余方程，并由正交幂等元识别环的直积分解。
\end{itemize}
\end{learninggoals}

在整数中，\(2\cdot3=6\) 不是零；把计算放到模 \(6\) 的剩余类里，它却变成零。这并不妨碍剩余类继续相加、相乘，但它说明从整数继承运算之后，乘法的性质可能改变。若要系统地问哪些元素能取逆、哪些乘积会成为零，以及哪些同余关系允许同时做加法和乘法，就需要把两种运算放进同一个结构中。理想控制可用于取商的同余关系，中国剩余定理则说明某些模数上的计算可以怎样拆开并重新合成。

\ProseParagraphBreak

本章使用第一编的交换群、同态、商群与整数子群的分类；矩阵和线性映射只需基本线性代数。环的乘法起初可以不交换，进入交换环的结论时会另行说明。一般极大理想的存在性证明使用 Zorn 引理\footnote{佐恩（Max August Zorn，1906-1993），德国数学家，研究代数并提出与选择公理等价的极大元引理。}。涉及准素理想的习题使用\cref{b1:def:primary-ideal}。

\ProseParagraphBreak
环的定义、基本例子与同态可对照席南华《基础代数》第一卷\cite[第5章第5.3节]{Xi2016BasicAlgebraI}。交换环中的理想、素理想、极大理想与商构造，可继续读 Atiyah 与 Macdonald 的《Introduction to Commutative Algebra》\cite[Chapter 1]{AtiyahMacdonald1969CommutativeAlgebra}；中国剩余定理的整数计算则可对照 Shoup 的《A Computational Introduction to Number Theory and Algebra》\cite[Section 4.3]{Shoup2008NumberTheoryAlgebra}。

% !TeX root = ../../Book1.tex
% 纲要标题：### 10.1 环的约定与例子
% 纲要位置：计划纲要.md，第 388 行。

\section{环的约定与例子}
\label{b1:sec:1001}

群只规定一种运算。整数与矩阵却同时允许加法和乘法，而且乘法能够分配到加法中。环的定义把这两种运算放在一起；它也使零因子、方程的可解性和同余关系进入同一套语言。本节先固定单位元及同态的约定，再比较几类常见例子。

% 纲要内容（仅作写作计划，尚非教材正文）：
%
% * 含幺环、交换环及同态的统一约定
% * 矩阵环、函数环、有限直积环、多项式环与群环
% * 零因子、单位与幂零元
%

\subsection{含幺环、交换环与同态}
\label{b1:subsec:1001:01}
\begin{definition}\label{b1:def:ring}
给定集合 \(R\) 及加法、乘法两个从 \(R\times R\) 到 \(R\) 的运算。如果它们满足下列条件，就称 \(R\) 连同这两种运算为\term[huan]{环}[ring]：
\begin{enumerate}
\item 加法结合且交换，存在零元 \(0\)，每个 \(a\in R\) 有加法逆元 \(-a\)；
\item 乘法结合，存在单位元 \(1\)，使 \(1a=a1=a\)；
\item 对任意 \(a,b,c\in R\)，有 \(a(b+c)=ab+ac\)、\((a+b)c=ac+bc\)。
\end{enumerate}
如果还对任意 \(a,b\in R\) 有 \(ab=ba\)，就称 \(R\) 为\term[jiaohuan-huan]{交换环}[commutative ring]。本书环总含幺，但允许 \(1=0\)；这时 \(R\) 只有一个元素，称为\term[linghuan]{零环}[zero ring]。
\end{definition}
加法公理正是说 \((R,+)\) 为交换群。由 \(a0=a(0+0)=a0+a0\)，在加法群中消去 \(a0\)，得到 \(a0=0\)；同样 \(0a=0\)。再由 \(a(b-b)=0\)，得 \(a(-b)=-(ab)\)，继而 \((-a)(-b)=ab\)。这些等式来自分配律，而不是额外公理。若 \(1=0\)，每个元素 \(a=a1=a0=0\)，说明零环的确只有一个元素。

\ProseParagraphBreak
允许零环会影响后面的边界条件：下面定义的整环、除环、域都要求非零，即 \(1\ne0\)。\secref{b1:sec:1002}将定义理想与真理想；零环的唯一理想就是整个环，因而没有真理想，极大理想的存在性也须在非零环中讨论。

\begin{definition}\label{b1:def:ring-homomorphism}
设 \(R,S\) 为环，\(f:R\rightarrow S\) 为映射。如果对任意 \(a,b\in R\) 都有
\[
f(a+b)=f(a)+f(b),\qquad f(ab)=f(a)f(b),\qquad f(1_R)=1_S,
\]
就称 \(f\) 为从 \(R\) 到 \(S\) 的\term[huan-tongtai]{环同态}[ring homomorphism]。如果环同态 \(f\) 还是双射，就称它为\term[huan-tonggou]{环同构}[ring isomorphism]。如果子集 \(T\subseteq R\) 含 \(1_R\)，并且对减法和乘法封闭，就称 \(T\) 为本书约定下的\term[zihuan]{子环}[subring]。
\end{definition}
同态保持零元和负元，因为它首先是加法群同态。双射环同态的逆映射也保持两种运算与单位：将等式写在原像上再应用 \(f\)，便逐项得到这些性质。有些文献允许环同态不保持单位；本书始终要求 \(f(1_R)=1_S\)，子环也须含母环的同一个单位。按这一约定，任意环到零环的唯一映射都是环同态，而零环到非零环没有环同态，因为它的同一个元素 \(0=1\) 不可能同时映到不同的 \(0\) 与 \(1\)。

\ProseParagraphBreak
要求保持单位使 \(\mathbb Z\rightarrow R\)、\(n\mapsto n1_R\) 成为唯一的环同态；这里正整数倍由重复加法定义，负数倍取其加法逆元。它说明环中的整数从何而来：像由全部 \(n1_R\) 组成，对减法和乘法封闭，并且每个含 \(1_R\) 的子环都必须包含它，所以它是由单位生成的最小子环。

\begin{definition}\label{b1:def:ring-characteristic}
设 \(R\) 为环。若存在正整数 \(n\) 使 \(n1_R=0\)，其中最小者称为 \(R\) 的\term[tezheng]{特征}[characteristic]；若不存在，规定特征为零。环 \(R\) 的特征记为 \(\operatorname{char}R\)。
\end{definition}
\symidx{\(\operatorname{char}R\)：环的特征}
同态 \(\mathbb Z\rightarrow R\) 的核是加法子群，\cref{b1:lem:integer-subgroups}说明它唯一等于 \(n\mathbb Z\)，其中 \(n\geq0\)；这与特征定义中的 \(n\) 相同。特征为零时，这个同态单射，所以由单位生成的子环与 \(\mathbb Z\) 同构；特征为正整数 \(n\) 时，对整数 \(a,b\) 有 \(a1_R=b1_R\) 当且仅当 \(a-b\in n\mathbb Z\)，故 \([a]_n\mapsto a1_R\) 给出从整数剩余类环 \(\mathbb Z/n\mathbb Z\) 到该子环的同构。它保持乘法是因为 \((a1_R)(b1_R)=(ab)1_R\)。特征不是集合的元素个数：例如 \(\mathbb Z\) 和 \(\mathbb Q\) 都具有特征零。零环的特征为一，对应 \(\mathbb Z/\mathbb Z\)。

\begin{definition}\label{b1:def:domain-field}
如果非零交换环 \(R\) 对任意 \(a,b\in R\) 都满足 \(ab=0\) 必有 \(a=0\) 或 \(b=0\)，就称 \(R\) 为\term[zhenghuan]{整环}[integral domain]。如果非零环 \(R\) 的每个非零元素都有双边乘法逆元，就称 \(R\) 为\term[chuhuan]{除环}[division ring]；乘法交换的除环称为\term[yu]{域}[field]。
\end{definition}
域是整环：若 \(a\ne0\) 且 \(ab=0\)，左乘 \(a^{-1}\) 得 \(b=0\)。\(\mathbb Z\) 是整环而不是域，\(\mathbb Q,\mathbb R,\mathbb C\) 是域。整环的特征只能为零或素数：若正特征 \(n=uv\)，\(1<u,v<n\)，则 \((u1_R)(v1_R)=0\)，而由正特征的最小性，两个因子都非零，与整环无零因子矛盾。

\ProseParagraphBreak
这几类环的条件可以按下面的蕴含关系比较：
\[
\begin{tikzcd}[column sep=2em,row sep=1.8em]
\text{域}\arrow[r,Rightarrow]\arrow[d,Rightarrow]
&\text{整环}\arrow[r,Rightarrow]
&\text{交换非零环}\arrow[d,Rightarrow]\\
\text{除环}\arrow[rr,Rightarrow]&&\text{非零环}
\end{tikzcd}
\]
域恰好是乘法交换的除环。整环只排除两个非零元素相乘为零的情形，不要求非零元素都可逆；除环要求每个非零元素都有双边逆，但不要求乘法交换。

\subsection{矩阵环、函数环、多项式环与群环}
\label{b1:subsec:1001:02}
设 \(R\) 为环，\(n\geq1\) 为整数。所有 \(n\times n\) 矩阵在逐项加法与通常的矩阵乘法下组成\term[juzhen-huan]{矩阵环}[matrix ring] \(M_n(R)\)，单位为单位矩阵。乘法结合律可由有限和核验：两种结合方式的第 \((i,j)\) 项都是 \(\sum_{k,\ell}a_{ik}b_{k\ell}c_{\ell j}\)。若 \(R\ne0,n\geq2\)，则矩阵单位满足 \(E_{12}E_{21}=E_{11}\)、\(E_{21}E_{12}=E_{22}\)，故即使 \(R\) 交换，矩阵环也不交换。这里 \(n\ge2\) 不可省略：取唯一矩阵项给出 \(M_1(R)\cong R\)，所以 \(1\times1\) 矩阵环保留系数环原来的交换性。
\symidx{\(M_n(R)\)：环上的 \(n\) 阶矩阵环}

\ProseParagraphBreak
设 \(R\) 为环，\(X\) 为集合。映射集合 \(R^X\) 在逐点加法和乘法下构成\term[hanshu-huan]{函数环}[function ring]，其零元与单位分别为常值函数 \(0,1\)。每条环公理都可在每个 \(x\in X\) 处验证。若 \(R\) 为域而 \(X\) 有两个不同的点 \(x,y\)，定义
\[
f(z)=\begin{cases}1,&z=x,\\0,&z\ne x\end{cases},
\qquad
g(z)=\begin{cases}1,&z=y,\\0,&z\ne y\end{cases}.
\]
两函数分别只在一个点非零，所以 \(f\ne0\)、\(g\ne0\)；因为 \(x\ne y\)，每个点处至少有一个函数取零，故 \(fg=0\)。因此，系数环是域也不能保证函数环是整环。

\ProseParagraphBreak
设 \(R_1,\ldots,R_r\) 为环，其中 \(r\geq1\)。在笛卡尔积 \(\prod_{i=1}^rR_i\) 上逐坐标定义加法与乘法：
\[
(a_i)_i+(b_i)_i=(a_i+b_i)_i,\qquad
(a_i)_i(b_i)_i=(a_ib_i)_i.
\]
零元和单位分别是 \((0_{R_i})_i\) 与 \((1_{R_i})_i\)。环公理逐坐标成立，所得环称为 \(R_1,\ldots,R_r\) 的\term[huan-zhiji]{直积环}[direct product of rings]；对每个 \(1\leq j\leq r\)，坐标投影 \(\prod_iR_i\to R_j\) 都是保单位的环同态。前述函数环 \(R^X\) 则是各坐标均取 \(R\)、指标集可任意选取的情形。

\ProseParagraphBreak
本章关于多项式只需建立形式和的运算、次数以及整环上的次数加法。除法算法、重根与不可约性将在\chapref{b1:ch:12}系统讨论；这里先把多项式环作为一种可以直接核验的环来构造。

\begin{definition}\label{b1:def:polynomial-ring}
设 \(R\) 为环，\(x\) 为与 \(R\) 中每个元素交换的形式符号。考虑系数 \(a_i\in R\) 且只有有限项非零的形式和
\(f=\sum_{i\geq0}a_ix^i\)，并规定两个形式和相等当且仅当对应系数逐项相等。再规定运算
\[
\sum_i a_ix^i+\sum_i b_ix^i=\sum_i(a_i+b_i)x^i,
\qquad
\Bigl(\sum_i a_ix^i\Bigr)\Bigl(\sum_j b_jx^j\Bigr)
=\sum_k\Bigl(\sum_{i+j=k}a_ib_j\Bigr)x^k.
\]
则称这些形式和组成的环为 \(R\) 上的\term[duoxiangshi-huan]{多项式环}[polynomial ring]，记为 \(R[x]\)。如果 \(R\) 不交换，乘积中的系数仍按所示次序相乘。
\end{definition}
\symidx{\(R[x]\)：一元多项式环}
这里的 \(x\) 是与全部系数交换的中心变量；这只规定 \(ax=xa\)，没有规定任意系数 \(a,b\) 也交换。例如
\[
(ax)b=a(xb)=a(bx)=(ab)x,
\]
所得系数仍按 \(a,b\) 的顺序相乘，不能改为 \(ba\)。

\ProseParagraphBreak
系数支集有限使上述各和有限。两次乘法的第 \(k\) 个系数无论如何加括号，都等于 \(\sum_{i+j+\ell=k}a_ib_jc_\ell\)，故乘法结合；分配律来自系数环，常数多项式 \(1\) 为单位。这验证了定义确实给出环。设 \(R\) 为环，\(f=\sum_{i\geq0}a_ix^i\in R[x]\) 为非零多项式，并令 \(n\) 为满足 \(a_n\ne0\) 的最大指标。整数 \(n\) 称为 \(f\) 的\term[duoxiangshi-cishu]{次数}[degree]，记为 \(\deg f\)；系数 \(a_n\) 称为 \(f\) 的\term[shouxiang-xishu]{首项系数}[leading coefficient]。如果 \(a_n=1\)，就称 \(f\) 为\term[shouyi-duoxiangshi]{首一多项式}[monic polynomial]。零多项式的次数记为 \(-\infty\)。
\symidx{\(\deg f\)：多项式的次数}

\begin{proposition}\label{b1:prop:polynomial-degree}
若 \(R\) 是整环，\(f,g\in R[x]\) 非零，则 \(fg\ne0\)，且
\(\deg(fg)=\deg f+\deg g\)。特别地，\(R[x]\) 是整环。
\end{proposition}
\begin{proof}
设 \(f,g\) 的次数为 \(m,n\)，首项系数为 \(a_m,b_n\)。乘积没有次数超过 \(m+n\) 的项，\(x^{m+n}\) 的系数是非零的 \(a_mb_n\)，所以最高次数恰为 \(m+n\)。
\end{proof}
形式多项式与多项式函数必须区别。例如在模 \(2\) 的剩余类中，\(x^2-x\) 对每个输入都取零值，却不是零多项式。若底环存在零因子，次数公式也可能失效：在模 \(4\) 的系数中，非零多项式 \(2x\) 与自身的乘积 \((2x)(2x)=4x^2\) 为零。

\ProseParagraphBreak
设 \(R\) 是交换环，\(G\) 为群。对系数 \(a_g\in R\)，若只有有限多个 \(a_g\) 非零，则称 \(\sum_{g\in G}a_gg\) 为有限形式和。全体这类形式和在逐项加法及
\[
\Bigl(\sum_g a_gg\Bigr)\Bigl(\sum_h b_hh\Bigr)
=\sum_{g,h}a_gb_h(gh)
\]
下组成\term[qunhuan]{群环}[group ring] \(R[G]\)。写在括号中的 \(gh\) 是群乘积，相同群元素的系数须相加。有限支集保证乘积有限，结合律由系数乘法与群乘法的结合律给出；单位是群单位处系数为 \(1\)、其余系数为零的形式和。这一构造把群的乘法保留下来，又允许对群元素作线性组合。例如取 \(G=C_2=\{1,g\}\)、\(g^2=1\)，对任意 \(a,b,c,d\in R\) 有
\[
(a+bg)(c+dg)
=ac+adg+bcg+bdg^2
=(ac+bd)+(ad+bc)g.
\]
群关系 \(g^2=1\) 把最后一项并入常数项。取 \(a=c=d=1\)、\(b=-1\)，便得到 \((1-g)(1+g)=0\)。
\symidx{\(R[G]\)：群环}

\subsection{零因子、单位与幂零元}
\label{b1:subsec:1001:03}
\begin{definition}\label{b1:def:units-zero-nil}
设 \(a\) 为环 \(R\) 的元素。如果存在 \(b\in R\) 使 \(ab=ba=1\)，就称 \(a\) 为\term[danwei]{单位}[unit]，所有单位组成的集合记为 \(R^\times\)。如果 \(a\ne0\)，且存在非零元素 \(b\in R\) 使 \(ab=0\) 或 \(ba=0\)，就称 \(a\) 为\term[lingyinzi]{零因子}[zero divisor]；如果存在正整数 \(m\) 使 \(a^m=0\)，就称 \(a\) 为\term[milingyuan]{幂零元}[nilpotent element]。
\end{definition}
\symidx{\(R^\times\)：环的单位群}
非交换环中，\(ab=0\) 与 \(ba=0\) 不能互相推出，文献中也会分别讨论左零因子与右零因子。本书在这里把这两种情形合称为零因子，只要求存在一个非零元素在至少一侧将它零化。

\ProseParagraphBreak
单位在乘法下组成群：乘积 \(ab\) 的逆是 \(b^{-1}a^{-1}\)，其余群公理从环继承。单位不可能为零因子，因为可以乘逆元消去。非零幂零元则一定是零因子：取使 \(a^m=0\) 的最小正整数 \(m\)，此时 \(m\geq2\)，最小性保证 \(a^{m-1}\ne0\)，而 \(a\cdot a^{m-1}=a^m=0\)。因此确实找到了一个非零的零化因子。注意零本身是幂零元，却按本书定义不叫零因子。

\ProseParagraphBreak
对照 Atiyah 与 Macdonald 的《Introduction to Commutative Algebra》时还须留意：该书在非零环中把零也计为零因子\cite[Chapter 1, p.~2]{AtiyahMacdonald1969CommutativeAlgebra}；本书的“零因子”始终要求元素非零。

\ProseParagraphBreak
有限性会把消去性质转化为可逆性，但非交换情形还须检查得到的是双边逆元。

\begin{proposition}\label{b1:prop:finite-domain-field}
有限非零环中，不是零因子的非零元素一定是单位。特别地，每个有限整环都是域。
\end{proposition}
\begin{proof}
设 \(a\ne0\) 且不是零因子。映射 \(L_a:r\mapsto ar\) 单射，因为 \(ar=as\) 给出 \(a(r-s)=0\)，从而 \(r=s\)。有限集上的单射为满射，故存在 \(b\) 使 \(ab=1\)。再比较
\[
L_a(ba)=a(ba)=(ab)a=a=L_a(1).
\]
由 \(L_a\) 的单射性，\(ba=1\)。这一步只用乘法结合律与单射性，没有使用交换性。整环乘法交换且每个非零元素都满足上述条件，故为域。
\end{proof}
有限性在这里不能删去：整数 \(2\) 既不是单位，也不是零因子。一般的非交换环中，仅有一侧逆元也不够。取域 \(k\) 上以 \(e_0,e_1,\ldots\) 为基的向量空间，令 \(T(e_i)=e_{i+1}\)，\(U(e_0)=0\)、\(U(e_{i+1})=e_i\)。在线性自映射组成的环中，\(UT=1\) 而 \(TU(e_0)=0\ne e_0\)，所以 \(TU\ne1\)。这说明定义中要求双边逆元确有必要。

\ProseParagraphBreak
幂零元还能产生单位。若 \(a^m=0\)，则
\[
(1-a)(1+a+\cdots+a^{m-1})
=(1+a+\cdots+a^{m-1})(1-a)=1.
\]
这里只涉及同一元素的幂，不需要环交换。与之相比，两个幂零元的和在非交换环中未必幂零。例如在任意域 \(k\) 上的 \(M_2(k)\) 中，取 \(a=E_{12}\)、\(b=E_{21}\)，则 \(a^2=b^2=0\)，而
\[
(a+b)^2=E_{11}+E_{22}=I_2.
\]
所以 \(a+b\) 不幂零。交换性怎样保证幂零条件对加法稳定，将在根理想的讨论中出现。

% !TeX root = ../../Book1.tex
% 纲要标题：### 10.2 理想与商环
% 纲要位置：计划纲要.md，第 394 行。

\section{理想与商环}
\label{b1:sec:1002}

在加法群上按一个子群取商，只能保证加法良定义。若还想让乘法通过商映射保留下来，这个子群必须能够吸收环元素的左右乘法，理想由此出现。商环中的零因子与可逆性又将分别引出素理想和极大理想。

% 纲要内容（仅作写作计划，尚非教材正文）：
%
% * 左理想、右理想与双边理想
% * 商环及同构定理
% * 素理想与极大理想的首次出现
%

\subsection{左理想、右理想与双边理想}
\label{b1:subsec:1002:01}
设 \(R\) 为环，先取加法子群 \(I\leq(R,+)\)，希望陪集上的乘法由代表元的乘法给出。若把第一项的代表 \(a\) 换成 \(a+\ell\)，其中 \(\ell\in I\)，所得乘积与原乘积的差是
\[
(a+\ell)b-ab=\ell b.
\]
要使输出陪集不变，必须有 \(\ell b\in I\)。若换的是第二项的代表，同样由 \(a(b+\ell)-ab=a\ell\) 得到 \(a\ell\in I\)。所以 \(I\) 既要吸收右乘，也要吸收左乘。与\cref{b1:thm:quotient-group}中由换代表元引出正规性一样，这里的条件也来自商运算的良定义要求。

\begin{definition}\label{b1:def:ring-ideal}
设 \(R\) 为环，\(I\subseteq R\) 是包含零元且对减法封闭的子集。若对每个 \(r\in R,a\in I\) 都有 \(ra\in I\)，称 \(I\) 为\term[zuo-lixiang]{左理想}[left ideal]；若总有 \(ar\in I\)，称为\term[you-lixiang]{右理想}[right ideal]；同时满足两者时称为\term[shuangbian-lixiang]{双边理想}[two-sided ideal]。未加限定的“理想”均指双边理想。若 \(I\ne R\)，称为真理想。
\end{definition}
理想不要求含环的单位；事实上理想一旦含有单位元或任一单位 \(u\)，就包含 \(u^{-1}u=1\)，继而包含每个 \(r=r1\)，所以等于 \(R\)。因此，一个非零环的真理想不是本书约定下的含同一单位的子环。这个限定不排除它自身的运算：理想仍对加法、减法和乘法封闭，可以作为不要求单位元的环来研究，有时也具有自己的单位元；这里排除的是它含有母环单位 \(1_R\)。

\ProseParagraphBreak
在交换环中左右条件等价。\(n\mathbb Z\) 是 \(\mathbb Z\) 的理想；反过来，整数环的每个理想首先是加法子群，由\cref{b1:lem:integer-subgroups}都为这种形式。

\ProseParagraphBreak
非交换情形确须区分左右。设 \(k\) 为域，在 \(M_2(k)\) 中令
\[
I_\ell=\left\{\,\begin{pmatrix}a&0\\b&0\end{pmatrix}
\;\middle|\;a,b\in k\,\right\},\qquad
I_r=\left\{\,\begin{pmatrix}a&b\\0&0\end{pmatrix}
\;\middle|\;a,b\in k\,\right\}.
\]
它们分别由第二列为零、第二行为零的矩阵组成，都对减法封闭。对任意 \(a,b,r,s,t,u\in k\)，
\[
\begin{aligned}
\begin{pmatrix}r&s\\t&u\end{pmatrix}
\begin{pmatrix}a&0\\b&0\end{pmatrix}
&=\begin{pmatrix}ra+sb&0\\ta+ub&0\end{pmatrix},\\
\begin{pmatrix}a&b\\0&0\end{pmatrix}
\begin{pmatrix}r&s\\t&u\end{pmatrix}
&=\begin{pmatrix}ar+bt&as+bu\\0&0\end{pmatrix}.
\end{aligned}
\]
第一式说明 \(I_\ell\) 吸收左乘，第二式说明 \(I_r\) 吸收右乘。但 \(E_{11}\in I_\ell\) 而 \(E_{11}E_{12}=E_{12}\notin I_\ell\)，所以 \(I_\ell\) 不是右理想；\(E_{11}\in I_r\) 而 \(E_{21}E_{11}=E_{21}\notin I_r\)，所以 \(I_r\) 不是左理想。

\ProseParagraphBreak
若 \(f:R\rightarrow S\) 为环同态，其加法核 \(\ker f=\{a\in R:f(a)=0\}\) 为双边理想：对 \(a\in\ker f\)，有 \(f(ra)=f(r)f(a)=0=f(ar)\)。这里保持零元的条件由加法保证，吸收乘法的条件来自乘法保持性。

\subsection{商环及同构定理}
\label{b1:subsec:1002:02}
\begin{proposition}\label{b1:prop:quotient-ring}
设 \(R\) 为环，\(L\leq(R,+)\) 为加法子群。在加法陪集集合 \(R/L\) 上按
\[
(a+L)+(b+L)=(a+b)+L,\qquad
(a+L)(b+L)=ab+L
\]
规定运算，则乘法与代表元无关当且仅当 \(L\) 是 \(R\) 的双边理想。在此情形下，\(R/L\) 构成\term[shanghuan]{商环}[quotient ring]，单位为 \(1+L\)；自然映射 \(q:R\rightarrow R/L\) 是满环同态，核为 \(L\)。
\end{proposition}
\begin{proof}
加法良定义来自加法群的商构造。若 \(L\) 是双边理想，而 \(a'=a+\ell,b'=b+m\)，其中 \(\ell,m\in L\)，则
\[
a'b'-ab=am+\ell b+\ell m\in L.
\]
其中 \(am\in L\) 用左吸收，因为 \(m\in L\)、\(a\in R\)；\(\ell b\in L\) 用右吸收，因为 \(\ell\in L\)、\(b\in R\)；\(\ell m\) 则用任一侧吸收都可得到。故乘法与代表元无关。结合律和分配律由代表元上的对应等式下降得到，\(1+L\) 在两侧乘任一陪集都保持它。自然映射按定义保持两种运算与单位，陪集都有代表元，且 \(a+L=L\) 恰好等价于 \(a\in L\)。

反过来，设所给陪集乘法良定义。对任意 \(r\in R\) 与 \(\ell\in L\)，由 \(\ell+L=0+L\) 得
\[
\begin{aligned}
r\ell+L&=(r+L)(\ell+L)=(r+L)(0+L)=L,\\
\ell r+L&=(\ell+L)(r+L)=(0+L)(r+L)=L.
\end{aligned}
\]
于是 \(r\ell,\ell r\in L\)。\(L\) 已是加法子群，故为双边理想。
\end{proof}
\symidx{\(R/I\)：环关于双边理想的商环}

把环同态 \(f:R\to S\) 先看作加法群同态，\cref{b1:thm:first-isomorphism}已经给出加法群同构 \((R,+)/\ker f\cong\operatorname{im}f\)。环论需要补的检查是：核吸收左右乘法，使商乘法有意义；诱导映射除保持加法外，还保持乘法和单位。前面已经证明核是双边理想，像则对两种运算封闭并含 \(1_S=f(1_R)\)，是 \(S\) 的子环。因而可以沿用同一个陪集映射得到下面的环同构。

\begin{theorem}[环的第一同构定理]\label{b1:thm:ring-first-iso}
设 \(R,S\) 为环。若 \(f:R\rightarrow S\) 是环同态，则
\[
\overline f:R/\ker f\longrightarrow\operatorname{im}f,
\qquad a+\ker f\longmapsto f(a)
\]
为环同构。更一般地，若 \(I\) 为 \(R\) 的双边理想且 \(I\subseteq\ker f\)，并以 \(q:R\rightarrow R/I\) 表示自然映射，则存在唯一环同态 \(g:R/I\rightarrow S\) 满足 \(g\circ q=f\)。
\end{theorem}
\begin{proof}
若 \(a-a'\in I\subseteq\ker f\)，则 \(f(a)=f(a')\)，所以 \(g(a+I)=f(a)\) 良定义。加法相容性由加法群的商构造给出，乘法与单位则分别满足
\[
\begin{aligned}
g\bigl((a+I)(b+I)\bigr)&=f(ab)=f(a)f(b),\\
g(1_R+I)&=f(1_R)=1_S.
\end{aligned}
\]
所以 \(g\) 是保单位的环同态。由定义有 \(g\circ q=f\)；由于 \(q\) 满射，该公式也是满足此等式的同态的唯一可能取值。取 \(I=\ker f\)，并将 \(g\) 的目标限制为其实际像 \(\operatorname{im}f\)，便得到定理中的 \(\overline f\)。它按像的定义满射；若 \(\overline f(a+\ker f)=0\)，则 \(a\in\ker f\)，所以其核只有零陪集，因而 \(\overline f\) 也是单射。
\end{proof}
条件 \(I\subseteq\ker f\) 使下面的三角形交换，虚线表示唯一满足 \(g\circ q=f\) 的环同态。
\[
\begin{tikzcd}[column sep=3em,row sep=2.5em]
R\arrow[r,"q"]\arrow[dr,"f"']&R/I\arrow[d,dashed,"\exists!\,g"]\\
&S
\end{tikzcd}
\]
这里的泛性质告诉我们：要从商环定义映射，先在原环定义，并检查需要抹去的理想确实映成零。商环不只是一个同构类型，还连同指定的商映射一起发挥作用。

\begin{proposition}[理想对应与商的同构]\label{b1:prop:ring-ideal-correspondence}
设 \(R\) 为环，\(I\) 为 \(R\) 的双边理想，\(q:R\rightarrow R/I\) 为自然映射。则：
\begin{enumerate}
\item 理想对应：\(q\) 给出包含 \(I\) 的 \(R\) 的双边理想与 \(R/I\) 的双边理想之间的保序双射。令 \(J\) 遍历前一类理想，\(L\) 遍历后一类理想，双向公式为
\begin{equation}\label{b1:eq:ring-ideal-correspondence}
J\longmapsto J/I,\qquad L\longmapsto q^{-1}(L).
\end{equation}
\item 环的第三同构定理：若 \(J\) 为 \(R\) 的双边理想且 \(I\subseteq J\)，则
\begin{equation}\label{b1:eq:ring-third-iso}
(R/I)/(J/I)\cong R/J.
\end{equation}
\item 环的第二同构定理：若 \(A\subseteq R\) 为含同一单位的子环，并记 \(A+I=\{a+i:a\in A,\ i\in I\}\)，则 \(A+I\) 是 \(R\) 的子环，\(I\) 是 \(A+I\) 的双边理想，且
\begin{equation}\label{b1:eq:ring-second-iso}
A/(A\cap I)\cong(A+I)/I.
\end{equation}
\end{enumerate}
\end{proposition}
\begin{proof}
\begin{enumerate}
\item 设 \(L\) 是 \(R/I\) 的双边理想。其原像 \(q^{-1}(L)\) 对减法及左右乘法封闭，且包含 \(I\)。含 \(I\) 的理想 \(J\) 的像 \(q(J)=J/I\) 对 \(R/I\) 中元素的左右乘法封闭，因为这些元素都可提升到 \(R\)。有 \(q^{-1}(q(J))=J\)：若 \(q(a)=q(j)\)，则 \(a-j\in I\subseteq J\)。又 \(q(q^{-1}(L))=L\)，因为 \(q\) 满射。两种操作均保持包含关系，得到\cref{b1:eq:ring-ideal-correspondence}。
\item 因 \(I\subseteq J\)，映射 \(R/I\rightarrow R/J\)、\(a+I\mapsto a+J\) 良定义且满射，核为 \(J/I\)。对该映射应用\cref{b1:thm:ring-first-iso}，即得\cref{b1:eq:ring-third-iso}。
\item 对 \(a,b\in A\) 及 \(i,j\in I\)，两元素 \(a+i,b+j\) 的差属于 \(A+I\)，乘积
\[
(a+i)(b+j)=ab+aj+ib+ij\in A+I.
\]
假设 \(A\) 含母环的同一个单位正用于这里：\(1_R\in A\subseteq A+I\)，配合上述封闭性，使 \(A+I\) 成为本书约定下的子环；只有加法与乘法封闭还不足以保证这一点。\(I\subseteq A+I\) 且吸收其左右乘法，因此是它的双边理想。\(A\cap I\) 同样是 \(A\) 的双边理想。映射 \(A\rightarrow(A+I)/I\)、\(a\mapsto a+I\) 保持单位且满射，因为 \(1_R\mapsto1_R+I\)、\(a+i+I=a+I\)，其核为 \(A\cap I\)。应用\cref{b1:thm:ring-first-iso}，即得\cref{b1:eq:ring-second-iso}。
\end{enumerate}
\end{proof}
上面的理想对应和两个商同构分别来自不同的映射。表中 \(I\subseteq J\) 是 \(R\) 的双边理想，\(A\) 是含 \(1_R\) 的子环；把出发映射及其核写出，就能看清每条公式抹去了什么。

\begin{center}
\begin{tabularx}{\linewidth}{@{}l>{\raggedright\arraybackslash}Xl@{}}
\toprule
结论 & 出发映射 & 核 \\
\midrule
理想对应 & \(q:R\to R/I\)，\(a\mapsto a+I\) & \(I\) \\
第三同构 & \(R/I\to R/J\)，\(a+I\mapsto a+J\) & \(J/I\) \\
第二同构 & \(A\to(A+I)/I\)，\(a\mapsto a+I\) & \(A\cap I\) \\
\bottomrule
\end{tabularx}
\end{center}

环同态的像未必是目标环的理想，例如 \(\mathbb Z\hookrightarrow\mathbb Q\) 的像不是 \(\mathbb Q\) 的理想。上述理想对应使用商映射的满射性，不能省去这个条件。

\subsection{素理想与极大理想}
\label{b1:subsec:1002:03}
本小节只讨论交换环。非交换环的素性与极大理想需要区分不同版本，不把下面的乘积判别直接当作其一般定义。

\ProseParagraphBreak
取真理想 \(I\) 后，商环 \(R/I\) 非零。接下来有两个不同的问题：什么时候它仍像整数环一样，两个非零元素的乘积不会为零？什么时候它像域一样，每个非零元素都能取逆？下面的素理想与极大理想将分别刻画这两种商环性质。

\begin{definition}\label{b1:def:prime-maximal-ideal}
设 \(R\) 为交换环，\(\mathfrak p\subsetneq R\) 为真理想。若对任意 \(a,b\in R\)，\(ab\in\mathfrak p\) 都蕴含 \(a\in\mathfrak p\) 或 \(b\in\mathfrak p\)，则称 \(\mathfrak p\) 为\term[su-lixiang]{素理想}[prime ideal]。再设 \(\mathfrak m\subsetneq R\) 为真理想。若满足 \(\mathfrak m\subseteq J\subseteq R\) 的理想只有 \(J=\mathfrak m\) 与 \(J=R\)，则称 \(\mathfrak m\) 为\term[jida-lixiang]{极大理想}[maximal ideal]。
\end{definition}
“真”不能漏掉；否则整个环会满足乘积条件，却不能对应非零整环。\((0)\) 是整环中的素理想，但整数环中它不极大，因为 \((0)\subsetneq2\mathbb Z\subsetneq\mathbb Z\)。

\begin{proposition}\label{b1:prop:prime-max-quotient}
交换环 \(R\) 的理想 \(I\) 为素理想，当且仅当 \(R/I\) 是整环；为极大理想，当且仅当 \(R/I\) 是域。因此每个极大理想都是素理想。
\end{proposition}
\begin{proof}
\(I\ne R\) 恰好保证商环非零；在商环中，\((a+I)(b+I)=0\) 恰好等价于 \(ab\in I\)，故第一项直接由定义得到。

若 \(I\) 极大且 \(a\notin I\)，要使非零类 \(a+I\) 有逆，就要找到 \(r\in R\) 使 \(ra-1\in I\)。这等价于 \(1\in I+Ra\)，所以考虑包含 \(I\) 并加入 \(a\) 的集合
\[
I+Ra=\{i+ra:i\in I,\ r\in R\}.
\]
此集合对减法封闭；由于 \(R\) 交换，对 \(s\in R\) 还有 \(s(i+ra)=si+(sr)a\in I+Ra\)，故它是理想。它含 \(I\) 及 \(a\)，而 \(a\notin I\)，所以严格包含 \(I\)，由极大性知 \(I+Ra=R\)。写 \(1=i+ra\)，便有 \((r+I)(a+I)=1+I\)，所以商中每个非零元素可逆。反过来，若 \(R/I\) 为域，而理想 \(J\) 严格包含 \(I\)，取 \(a\in J\setminus I\)；商中 \(a+I\) 的逆有代表 \(r\)，于是 \(1-ra\in I\subseteq J\)，给出 \(1\in J\)，故 \(J=R\)。最后，域无非零零因子，所以极大理想为素理想。
\end{proof}
素理想不一定极大，正如整环不一定是域：\(\mathbb Z/(0)\cong\mathbb Z\) 是整环，却不是域，所以 \((0)\) 是 \(\mathbb Z\) 的素理想而不是极大理想。这与前面的严格包含链 \((0)\subsetneq2\mathbb Z\subsetneq\mathbb Z\) 给出同一个反例。

\ProseParagraphBreak
整数环的理想为 \(n\mathbb Z\)。当 \(p\) 是素数，任一 \(a\notin p\mathbb Z\) 满足 \(\gcd(a,p)=1\)，整数 Bézout 等式给出 \(ua+vp=1\)，故模 \(p\) 的非零类均有逆。所得域记为 \(\mathbb F_p=\mathbb Z/p\mathbb Z\)。正合数 \(n=ab\)、\(1<a,b<n\) 则产生两个非零类之积为零。因此 \(\mathbb Z\) 的素理想恰为 \((0)\) 及 \(p\mathbb Z\)，极大理想恰为后者。
\symidx{\(\mathbb F_p\)：含 \(p\) 个元素的素数域}

\ProseParagraphBreak
极大理想的存在性寻找的是包含偏序中的极大元，即不能再严格扩大的真理想，并非包含所有真理想的最大元。例如 \(2\mathbb Z\) 与 \(3\mathbb Z\) 都是极大理想，却互不包含；没有哪一个同时包含这两个真理想并仍为真理想。下面用 Zorn 引理证明所需极大元存在。

\begin{proposition}\label{b1:prop:maximal-ideal-exists}
非零交换环中的每个真理想都包含在某个极大理想中。
\end{proposition}
\begin{proof}
固定真理想 \(I\)，考虑包含 \(I\) 的真理想按包含关系构成的偏序集。它非空。非空链的并对减法封闭，因为任意两个元素都落在链上某个共同的理想中；它也吸收乘法。这个并不含 \(1\)，否则链中的某个真理想已含 \(1\)，矛盾。因而链的并仍是此偏序集的上界；空链可用 \(I\) 作上界。应用 Zorn 引理（\cref{b1:thm:A-zorn}），即“每条链有上界的非空偏序集有极大元”，得到 \(\mathfrak m\)。任何严格包含它的真理想仍包含 \(I\)，违背极大性，故 \(\mathfrak m\) 是极大理想。
\end{proof}
这里明确使用选择公理的 Zorn 形式，不给出寻找极大理想的算法。在具体环中仍应尽可能写出商映射，例如对交换环 \(R\)，常数项映射 \(R[x]\rightarrow R\) 满射，核为 \(xR[x]\)。\cref{b1:thm:ring-first-iso}给出 \(R[x]/xR[x]\cong R\)，所以当 \(R\) 为整环或域时，这个理想分别为素理想或极大理想。

% !TeX root = ../../Book1.tex
% 纲要标题：### 10.3 理想运算
% 纲要位置：计划纲要.md，第 400 行。

\section{理想运算}
\label{b1:sec:1003}

一个理想记录一组允许被视为零的元素。设 \(I,J\) 为环 \(R\) 的双边理想。若在同一个商环中令 \(I\) 和 \(J\) 的元素都成为零，就要商去 \(I+J\)；同时属于 \(I\) 和 \(J\) 的元素则组成 \(I\cap J\)，它们在两个商环中都映为零。具体地，自然同态 \(r\mapsto(r+I,r+J)\) 满足
\[
\ker\!\left(R\longrightarrow R/I\times R/J\right)=I\cap J.
\]
理想积 \(IJ\) 由 \(I\) 中元素与 \(J\) 中元素的乘积的有限和组成。从包含关系看，\(I+J\) 是同时包含 \(I,J\) 的最小理想，\(I\cap J\) 是同时包含于 \(I,J\) 的最大理想；\(IJ\) 则由环的乘法产生，并且包含于 \(I\cap J\)。和、交由包含关系确定，积回答的则是乘法问题。本节先写出它们的元素形式，再在交换环中讨论取幂、取根，以及寻找哪些乘子能把一个理想送入另一个理想。

% 纲要内容（仅作写作计划，尚非教材正文）：
%
% * 和、积、交与生成理想
% * 理想的幂与根、准素理想的初步性质
% * 不同理想运算之间的包含关系
% * 理想商、零化子与基本运算性质
%

\subsection{理想的和、积、交与生成}
\label{b1:subsec:1003:01}
\begin{definition}\label{b1:def:ideal-operations}
设 \(I,J\) 为环 \(R\) 的双边理想。它们的\term[lixiang-de-he]{和}[sum of ideals]与\term[lixiang-de-ji]{积}[product of ideals]分别为
\[
I+J=\{\,a+b\mid a\in I,\ b\in J\,\},\qquad
IJ=\Bigl\{\,\sum_{\nu=1}^m a_\nu b_\nu\mid
m\geq0,\ a_\nu\in I,\ b_\nu\in J\,\Bigr\}.
\]
空和表示零。理想的交采用通常的集合交。对任意子集 \(S\subseteq R\)，包含 \(S\) 的所有双边理想之交称为 \(S\) 的\term[shengcheng-lixiang]{生成理想}[generated ideal]，记为 \((S)\)。
\end{definition}
\symidx{\(I+J,\ IJ,\ I\cap J\)：理想的和、积与交}
\symidx{\((S)\)：由子集 \(S\) 生成的理想}

生成双边理想时，两侧系数来自左右吸收条件：一个包含 \(s\in S\) 的双边理想必须包含所有 \(rs\) 和 \(st\)，继而包含所有 \(rst\)，其中 \(r,t\in R\)。再把这些元素的有限和收入其中，就得到下面的具体公式；非交换环中通常不能把两侧系数越过 \(s\) 合并。

\begin{proposition}\label{b1:prop:ideal-operations}
设 \(I,J\) 为环 \(R\) 的双边理想，\(S\subseteq R\)。则 \(I+J\)、\(IJ\) 与 \(I\cap J\) 都是 \(R\) 的双边理想；\(I+J\) 是包含 \(I,J\) 的最小双边理想。由 \(S\) 生成的双边理想有具体表达式
\[
(S)=\Bigl\{\,\sum_{\nu=1}^m r_\nu s_\nu t_\nu
\mid m\geq0,\ r_\nu,t_\nu\in R,\ s_\nu\in S\,\Bigr\}.
\]
若 \(R\) 交换，可把每个 \(r_\nu t_\nu\) 合成一个系数，得到 \(\sum r_\nu s_\nu\) 的形式。
\end{proposition}
\begin{proof}
理想的交保持减法及左右乘法封闭。对 \(I+J\)，两个和相减仍可分别归入 \(I,J\)，且左右乘法分配到两项中，故是理想；任何同时包含 \(I,J\) 的理想都包含这些和。

\(IJ\) 对加法封闭来自有限和的合并，取负可把负号吸收到 \(a_\nu\) 中。对任意 \(r\in R\)，左乘将 \(a_\nu\) 换成 \(ra_\nu\in I\)，右乘将 \(b_\nu\) 换成 \(b_\nu r\in J\)，所以也是理想。最后一个显示式右边同样对这些运算封闭，取 \(r_\nu=t_\nu=1\) 说明它包含 \(S\)；任何包含 \(S\) 的双边理想都包含这些有限和，故恰为 \((S)\)。
\end{proof}
即使 \(S\) 无限，每个生成理想中的元素也只用有限多个生成元。写 \((a_1,\ldots,a_m)\) 表示有限集合的生成理想；由一个元素生成的理想称为\term[zhulixiang]{主理想}[principal ideal]。交换环中 \((a)=Ra\)，非交换环中的双边主理想则一般须保留两侧系数。

\ProseParagraphBreak
理想积取有限和，是为了保证加法封闭。一个具体例子在 \(\mathbb Q[x,y]\) 中可以看清。这里取迭代多项式环 \((\mathbb Q[x])[y]\)，其中 \(y\) 与系数交换，每个元素都唯一写成有限和 \(\sum_{i,j}a_{ij}x^iy^j\)。令 \(I=J=(x,y)\)。单个乘积的集合含有 \(x^2=x\cdot x\)、\(y^2=y\cdot y\)，却不含它们的和 \(x^2+y^2\)。事实上，若 \(x^2+y^2=fg\)，其中 \(f,g\in(x,y)\)，则 \(f,g\) 的常数项都为零。记它们的一次项分别为 \(\alpha x+\beta y\)、\(\gamma x+\delta y\)，其中系数均为有理数。乘积的二次项只能来自这两组一次项相乘，比较系数便有
\[
\alpha\gamma=1,\qquad
\alpha\delta+\beta\gamma=0,\qquad
\beta\delta=1.
\]
首尾两式使四个系数都非零，且 \(\gamma=\dfrac1\alpha\)、\(\delta=\dfrac1\beta\)。代入中间一式并乘以 \(\alpha\beta\)，得到 \(\alpha^2+\beta^2=0\)，与非零有理数的平方为正矛盾。因此 \(x^2+y^2\) 属于 \(IJ\)，但不属于单个乘积的集合。

\ProseParagraphBreak
有限和的定义还使
\[
I(J+K)=IJ+IK,\qquad (I+J)K=IK+JK,\qquad (IJ)K=I(JK).
\]
前两式把一个乘积的因子展开，再合并有限和；末式两边都是元素 \(abc\)、\(a\in I,b\in J,c\in K\) 的有限和。这些论证只用结合律和分配律，不要求交换性。

\ProseParagraphBreak
再设 \(k\) 为任意域，仍取迭代多项式环 \(k[x,y]=(k[x])[y]\)。在这个环中取 \(I=(x,y)\)、\(J=(x,y^2)\)。把两组生成元两两相乘，得到
\[
IJ=(x^2,xy^2,xy,y^3)=(x^2,xy,y^3).
\]
其中 \(xy\) 是容易漏掉的混合乘积；\(xy^2\) 已由 \(xy\) 生成，因此可从生成元列表中删去。

\subsection{理想的幂与根}
\label{b1:subsec:1003:02}
设 \(I\) 为环 \(R\) 的双边理想。约定 \(I^0=R\)、\(I^1=I\)，并对整数 \(n\geq1\) 递归定义 \(I^{n+1}=I^nI\)。对 \(n\geq1\)，\(I^n\) 由 \(n\) 个 \(I\) 中元素的乘积的有限和组成，而不只由某个元素的 \(n\) 次幂组成。交换环的主理想满足 \((a)^n=(a^n)\)，但多生成元时各个混合乘积也必须计算。

\begin{definition}\label{b1:def:radical-ideal}
设 \(R\) 为交换环，\(I\) 为 \(R\) 的理想。集合
\[
\sqrt I\coloneqq\{\,a\in R\mid a^n\in I\;\text{对某个正整数}\;n\;\text{成立}\,\}
\]
称为 \(I\) 的\term[lixiang-de-gen]{根}[radical of an ideal]。若 \(I=\sqrt I\)，则称 \(I\) 为\term[genlixiang]{根理想}[radical ideal]。\(\sqrt{(0)}\) 称为环的\term[milinggen]{幂零根}[nilradical]。
\end{definition}
\symidx{\(\sqrt{(0)}\)：交换环的幂零根}
\symidx{\(\sqrt I\)：理想的根}

取根可以直接从商环来理解。设 \(q:R\to R/I\) 为商映射，则
\[
a\in\sqrt I
\quad\Longleftrightarrow\quad
(a+I)^n=0\;\text{对某个正整数}\;n\;\text{成立}.
\]
因此 \(\sqrt I\) 收集商环中所有幂零元的原像：一个元素未必已经属于 \(I\)，但它在商环中的像经过某次幂会变为零。

\begin{proposition}\label{b1:prop:radical-properties}
设 \(R\) 为交换环，\(I\) 为 \(R\) 的理想。则 \(\sqrt I\) 是包含 \(I\) 的理想，且 \(\sqrt{\sqrt I}=\sqrt I\)。它是包含 \(I\) 的最小根理想。商环 \(R/I\) 中没有非零幂零元，当且仅当 \(I=\sqrt I\)。
\end{proposition}
\begin{proof}
若 \(a^m,b^n\in I\)，则交换性允许对 \((a-b)^{m+n-1}\) 用二项式展开。指数 \(m+n-1\) 的选择保证：若某项中 \(a\) 的指数小于 \(m\)，则至多为 \(m-1\)，因而 \(b\) 的指数至少为 \((m+n-1)-(m-1)=n\)。所以每项或者含 \(a^m\)，或者含 \(b^n\)，都属于 \(I\)。又 \((ra)^m=r^ma^m\in I\)，所以 \(\sqrt I\) 对减法及乘法吸收封闭。\(I\subseteq\sqrt I\) 可取指数一。

若 \(a^r\in\sqrt I\)，再取 \(s\) 使 \((a^r)^s\in I\)，得到 \(a\in\sqrt I\)，证明取两次根不变。若根理想 \(J\) 包含 \(I\)，则 \(a^n\in I\) 蕴含 \(a^n\in J\)，进而 \(a\in J\)，故 \(\sqrt I\subseteq J\)。最后，\((a+I)^n=0\) 等价于 \(a^n\in I\)，而 \(a+I=0\) 等价于 \(a\in I\)，给出商环的判别。
\end{proof}
没有非零幂零元的交换环称为\term[yuehua-huan]{约化环}[reduced ring]。对交换环 \(R\) 及其理想 \(I\)，\cref{b1:prop:radical-properties}给出
\[
\begin{aligned}
R\;\text{约化}&\quad\Longleftrightarrow\quad(0)\;\text{是根理想},\\
R/I\;\text{约化}&\quad\Longleftrightarrow\quad I\;\text{是根理想}.
\end{aligned}
\]
这里 \(\sqrt I/I\) 正是 \(R/I\) 中所有幂零元组成的理想。整环必约化，反之不成立：域的直积 \(k\times k\) 有零因子 \((1,0),(0,1)\)，但若 \((a,b)^n=(0,0)\)，域中取幂不产生非零幂零元，故 \(a=b=0\)。

\ProseParagraphBreak
每个素理想都是根理想，因为 \(a^n\in\mathfrak p\) 可反复用素性推出 \(a\in\mathfrak p\)。\(\mathbb Z\) 中，\(\sqrt{12\mathbb Z}=6\mathbb Z\)：若 \(12\mid a^n\)，素数 \(2,3\) 均整除 \(a\)；反过来 \(6\mid a\) 时 \(12\mid a^2\)。这也说明根理想的作用是忘掉某些幂次信息，而不是忘掉全部整除信息。

\ProseParagraphBreak
若希望分解理想而保留幂次信息，就会用到下述概念。与素理想要求商环没有零因子相比，准素条件允许出现零因子，但要求它们经过某次幂变为零。先在原环中写出这个条件，再翻译到商环。
\begin{definition}\label{b1:def:primary-ideal}
设 \(Q\) 为交换环 \(R\) 的真理想。若对任意 \(a,b\in R\)，只要 \(ab\in Q\) 且 \(a\notin Q\)，就存在正整数 \(m\) 使 \(b^m\in Q\)，则称 \(Q\) 为\term[zhunsu-lixiang]{准素理想}[primary ideal]。
\end{definition}
设 \(R\) 为交换环，\(Q\subsetneq R\) 为理想。准素性的商环判别为
\[
Q\;\text{是准素理想}
\quad\Longleftrightarrow\quad
R/Q\;\text{中每个非零零因子都是幂零元}.
\]
若 \(Q\) 准素，取商环中的一个非零零因子 \(b+Q\)。存在非零元素 \(a+Q\) 使 \((a+Q)(b+Q)=0\)，即 \(ab\in Q\) 且 \(a\notin Q\)。准素性给出正整数 \(m\) 使 \(b^m\in Q\)，所以 \((b+Q)^m=0\)，这个零因子幂零。

\ProseParagraphBreak
反过来，若 \(R/Q\) 的每个非零零因子都幂零，设 \(ab\in Q\) 且 \(a\notin Q\)。若 \(b\in Q\)，取 \(m=1\) 就满足准素条件；若 \(b\notin Q\)，则 \(a+Q,b+Q\) 都非零，而它们的乘积为零，故 \(b+Q\) 是非零零因子。由假设存在正整数 \(m\) 使 \((b+Q)^m=0\)，也就是 \(b^m\in Q\)。因此 \(Q\) 是准素理想。两种情况都要保留，因为本书的零因子定义不把零本身计入其中。

\begin{proposition}\label{b1:prop:primary-radical}
设 \(R\) 为交换含幺环。\(R\) 的每个素理想都是准素理想；若 \(Q\) 是 \(R\) 的准素理想，则 \(\sqrt Q\) 是素理想。因此，准素理想 \(Q\) 是素理想，当且仅当它是根理想。
\end{proposition}
\begin{proof}
若 \(P\) 为素理想，\(ab\in P\) 且 \(a\notin P\) 时有 \(b\in P\)，所以 \(P\) 满足准素条件。设 \(Q\) 为准素理想。由 \(Q\neq R\) 可知 \(1\notin\sqrt Q\)，故 \(\sqrt Q\) 仍是真理想。若 \(ab\in\sqrt Q\) 而 \(a\notin\sqrt Q\)，取正整数 \(n\) 使 \(a^nb^n\in Q\)。此时 \(a^n\notin Q\)，准素性给出正整数 \(m\)，使 \(b^{nm}\in Q\)；于是 \(b\in\sqrt Q\)。这证明 \(\sqrt Q\) 为素理想。若 \(Q\) 还是根理想，则 \(Q=\sqrt Q\) 为素理想；反之，素理想是根理想。
\end{proof}
例如整数环中的 \(p^r\mathbb Z\)（\(p\) 为素数，\(r\geq1\)）是准素理想。若 \(p^r\mid ab\) 而 \(p^r\nmid a\)，整数的素数分解说明 \(p\mid b\)。写 \(b=pc\)，便有 \(b^r=p^rc^r\)，所以 \(p^r\mid b^r\)；通过取幂，就补足了原来可能不足的素数因子个数。特别地，
\[
12\mathbb Z=4\mathbb Z\cap3\mathbb Z
\]
把一个理想写成两个准素理想的交，仍保留因子二所带的平方信息。这里仅用这个例子说明分解问题；一般的准素分解理论，包括 Noether 环中的存在性与唯一性范围，安排在第三卷第4.2--4.4节。

\subsection{理想运算的包含关系}
\label{b1:subsec:1003:03}
\begin{proposition}\label{b1:prop:ideal-containments}
设 \(I,J\) 为环 \(R\) 的双边理想，\(n\geq1\) 为整数。则
\[
IJ\subseteq I\cap J\subseteq I+J,\qquad I^{n+1}\subseteq I^n.
\]
如果 \(R\) 交换，还有
\[
\sqrt{IJ}=\sqrt{I\cap J}=\sqrt I\cap\sqrt J,
\qquad \sqrt{I^n}=\sqrt I\quad(n\geq1).
\]
\end{proposition}
\begin{proof}
每个 \(ab\)、\(a\in I,b\in J\) 既在 \(I\) 中也在 \(J\) 中，因此有限和也在交中；交中的元素可写为自身加零，故在和中。把 \(J\) 换成 \(I\) 的幂给出递降关系。

由包含关系及根的定义，\(\sqrt{IJ}\subseteq\sqrt{I\cap J}\subseteq\sqrt I\cap\sqrt J\)。若 \(a^r\in I\)、\(a^s\in J\)，则 \(a^{r+s}=a^ra^s\in IJ\)，故反向包含也成立。\(I^n\subseteq I\) 给出一个方向；若 \(a^r\in I\)，则 \(a^{rn}\in I^n\)，给出另一个方向。
\end{proof}
\cref{b1:fig:ideal-operation-containments}把和、交、积的包含关系放在一起。箭头沿包含方向由较小的理想指向较大的理想，图中的包含一般都可能严格；和与交分别是共同的最小上界与最大下界，而积位于交之内。

\begin{figure}[H]
\centering
\begin{tikzpicture}[x=1cm,y=1cm,>=stealth]
\node (sum) at (0,3) {\(I+J\)};
\node (leftideal) at (-1.5,2) {\(I\)};
\node (rightideal) at (1.5,2) {\(J\)};
\node (intersection) at (0,1) {\(I\cap J\)};
\node (product) at (0,0) {\(IJ\)};
\draw[->] (product) -- (intersection);
\draw[->] (intersection) -- (leftideal);
\draw[->] (intersection) -- (rightideal);
\draw[->] (leftideal) -- (sum);
\draw[->] (rightideal) -- (sum);
\end{tikzpicture}
\caption[理想运算的包含关系]{理想运算的包含关系。\(I,J\) 为同一环的双边理想，箭头表示包含。}
\label{b1:fig:ideal-operation-containments}
\end{figure}

等号需要额外理由。在 \(\mathbb Z\) 中取 \(I=J=2\mathbb Z\)，则 \(IJ=4\mathbb Z\subsetneq I\cap J=2\mathbb Z\)；取 \(I=2\mathbb Z,J=3\mathbb Z\)，则交为 \(6\mathbb Z\)，和为 \(\mathbb Z\)，也不相等。另一方面 \(I\cup J\) 通常不是理想，例如上述并包含 \(2,3\)，却不包含 \(3-2=1\)。

\ProseParagraphBreak
主理想的包含关系对应于生成元之间的整除关系，但方向相反。整数环中 \((a)\subseteq(b)\) 恰好等价于 \(b\mid a\)。例如 \(6\mathbb Z\subseteq2\mathbb Z\)，不是因为 \(6\mid2\)，而是每个 \(6\) 的倍数都是 \(2\) 的倍数。下一章将把这个简单现象推广到整环。下面先从乘法吸收条件出发，再定义一种理想运算。

\subsection{理想商与零化子}
\label{b1:subsec:1003:04}

本小节始终设 \(R\) 为交换含幺环。给定理想 \(I,J\)，要寻找满足 \(KJ\subseteq I\) 的最大理想 \(K\)。逐元素看，就是考虑哪些 \(r\in R\) 能使 \(J\) 中每个元素乘以 \(r\) 后都进入 \(I\)。把所有这些乘子收集起来，就得到所求的最大理想；下面先给出集合公式，再验证它的理想性质与最大性。

\begin{definition}\label{b1:def:ideal-quotient}
设 \(R\) 为交换含幺环，\(I,J\) 为 \(R\) 的理想。集合
\[
(I:J)\coloneqq\{\,r\in R\mid rJ\subseteq I\,\}
\]
称为\term[lixiang-shang]{理想商}[ideal quotient]，也称商理想。对单个元素 \(a\in R\)，简记
\[
I:a\coloneqq\bigl(I:(a)\bigr)
=\{\,r\in R\mid ra\in I\,\}.
\]
取 \(I=(0)\)，所得理想称为 \(J\) 的\term[linghuazi]{零化子}[annihilator]，也称湮灭子，记为 \(\operatorname{Ann}_R(J)=(0:J)\)；元素的零化子相应记为 \(\operatorname{Ann}_R(a)=0:a\)。把乘以 \(r\) 后使 \(a\) 变为零称为 \(r\) 对 \(a\) 的\term[linghua]{零化}[annihilation]，条件就是 \(ra=0\)。
\end{definition}
\termidx[shang-lixiang]{商理想}
\termidx[yanmiezi]{湮灭子}
\symidx{\((I:J)\)：理想商}
\symidx{\(\operatorname{Ann}_R(J)\)：理想的零化子}
这里的 \((I:J)\) 是 \(R\) 的一个子集，不是商环 \(R/I\)，也不是对两个理想作加法群的商。它的定义不要求 \(I\subseteq J\) 或 \(J\subseteq I\)。单元素简记中的等号成立，是因为 \(ra\in I\) 蕴含对每个 \(s\in R\) 有 \(rsa\in I\)，而反向可取 \(s=1\)。

\begin{proposition}\label{b1:prop:ideal-quotient-rules}
设 \(I,J\) 为交换含幺环 \(R\) 的理想。则 \((I:J)\) 是包含 \(I\) 的理想。对 \(R\) 的任一理想 \(K\)，有
\begin{equation}\label{b1:eq:ideal-quotient-adjunction}
KJ\subseteq I\quad\Longleftrightarrow\quad K\subseteq(I:J).
\end{equation}
因此 \((I:J)\) 是满足左侧包含关系的最大理想，且
\[
(I:R)=I,\qquad (I:0)=R,\qquad
(I:J)=R\ \Longleftrightarrow\ J\subseteq I.
\]
\end{proposition}
\begin{proof}
零在 \((I:J)\) 中。若 \(r,r'\) 在其中，对任意 \(j\in J\) 有 \((r-r')j=rj-r'j\in I\)；若 \(s\in R\)，则 \((sr)j=s(rj)\in I\)。所以它是理想。若 \(r\in I\)，理想的吸收性给出 \(rJ\subseteq I\)，故 \(I\subseteq(I:J)\)。

条件 \(KJ\subseteq I\) 意味着每个 \(k\in K\) 都满足 \(kJ\subseteq I\)，故 \(k\in(I:J)\)。反过来，若所有这些乘积均在 \(I\) 中，其有限和也在 \(I\) 中，正好得到 \(KJ\subseteq I\)。取 \(K=I:J\) 及任一满足条件的 \(K\)，分别得到它满足该条件和它包含全部这样的理想。最后，\(rR\subseteq I\) 等价于 \(r\in I\)，因为可令乘数为一；乘以零没有限制；\(1\in(I:J)\) 则等价于 \(J\subseteq I\)。这证明全部公式。
\end{proof}
例如在 \(\mathbb Z\) 中，
\[
\bigl(12\mathbb Z:8\mathbb Z\bigr)=3\mathbb Z,
\]
因为 \(12\mid8r\) 当且仅当 \(3\mid r\)。这里计算的是满足 \(K(8\mathbb Z)\subseteq12\mathbb Z\) 的最大理想 \(K\)，普通数字除法的结果 \(\dfrac{12}{8}=\dfrac32\) 不参与这个乘子条件。事实上，\((3\mathbb Z)(8\mathbb Z)=24\mathbb Z\subsetneq12\mathbb Z\)，所以把理想商再乘回去，一般只能得到包含而非等号。在 \(\mathbb Z/12\mathbb Z\) 中，\(\operatorname{Ann}_R(\bar4)=(\bar3)\)：条件 \(4r\equiv0\pmod{12}\) 正好要求 \(3\mid r\)。对非零元素，零化子非零恰好说明它是零因子；零元素本身的零化子则是全环。

\begin{proposition}\label{b1:prop:ideal-quotient-identities}
设 \(I,J,K\) 为交换含幺环 \(R\) 的理想。则理想商满足
\[
\begin{aligned}
(I\cap J):K&=(I:K)\cap(J:K),\\
I:(J+K)&=(I:J)\cap(I:K),\\
(I:J):K&=I:(JK).
\end{aligned}
\]
它对第一项保持包含关系，对第二项反转包含关系。若 \(m\geq0\) 为整数，且 \(J=(a_1,\ldots,a_m)\)，则
\[
(I:J)=\bigcap_{i=1}^m(I:a_i).
\]
\end{proposition}
\begin{proof}
一个元素 \(r\) 属于第一式左边，当且仅当 \(rK\) 同时包含于 \(I,J\)，这正是右边的两个条件。第二式中，\(r(J+K)\subseteq I\) 蕴含 \(rJ,rK\subseteq I\)，因为 \(J,K\subseteq J+K\)；反向由加法封闭性得到。

第三式左边的条件是对所有 \(k\in K\) 都有 \(rk\in(I:J)\)，即对所有 \(k\in K,j\in J\) 都有 \(rkj\in I\)。这些乘积的有限和仍在 \(I\) 中，故等价于 \(r(JK)\subseteq I\)。包含方向则由定义可见：扩大允许落入的理想放宽条件，扩大被乘的理想增加条件。最后，\(rJ\subseteq I\) 必须使每个 \(ra_i\in I\)；若每个生成元都满足此条件，则 \(r\sum_i s_i a_i=\sum_i s_i ra_i\in I\)，所以也足够。空生成集对应零理想，此时空交取为 \(R\)，公式仍成立。
\end{proof}
由\cref{b1:prop:ideal-quotient-rules}，\(I\subseteq(I:a)\)。所以商映射下的像 \((I:a)/I=\{\,r+I\mid r\in(I:a)\,\}\) 是 \(R/I\) 中的理想，这也符合\cref{b1:prop:ring-ideal-correspondence}中的理想对应。在商环中，\(r+I\) 零化 \(a+I\) 恰好意味着 \(ra\in I\)，也就是 \(r\in(I:a)\)。因而
\[
\operatorname{Ann}_{R/I}(a+I)=(I:a)/I.
\]
这把理想商与商环中的零因子联系起来。到\secref{b1:sec:1304}，允许乘以任意高次的分母，还会把这些理想商连接成饱和理想；下一节先继续讨论互素理想与环的直积分解。

% !TeX root = ../../Book1.tex
% 纲要标题：### 10.4 中国剩余定理
% 纲要位置：计划纲要.md，第 406 行。

\section{中国剩余定理}
\label{b1:sec:1004}

整数的两个互素模数可以独立指定余数。对一般交换环，替代“模数互素”的条件是两个理想的和等于整个环。中国剩余定理不仅给出同余方程的解，还把一个商环拆成若干分量；逆向看，这些分量由特殊的幂等元识别。本节除最后的辨析外均在交换含幺环中讨论。

% 纲要内容（仅作写作计划，尚非教材正文）：
%
% * 两两互素理想与有限直积分解
% * 整数同余方程
% * 幂等元与环的分解
%
% ---
%

\subsection{两两互素理想与有限直积分解}
\label{b1:subsec:1004:01}
\begin{definition}\label{b1:def:comaximal-ideals}
设 \(R\) 为交换环，\(I,J\) 为 \(R\) 的理想。若 \(I+J=R\)，则称 \(I,J\) 为\term[husu-lixiang]{互素理想}[comaximal ideals]。这等价于存在 \(a\in I,b\in J\) 使 \(a+b=1\)。
\end{definition}

\begin{lemma}\label{b1:lem:comaximal-product}
设 \(I,J\) 为交换环 \(R\) 的互素理想，则 \(I\cap J=IJ\)。更一般地，若 \(r\) 为正整数，且 \(R\) 的理想 \(I_1,\ldots,I_r\) 两两互素，则
\(\bigcap_{i=1}^r I_i=\prod_{i=1}^r I_i\)。
\end{lemma}
\begin{proof}
\cref{b1:prop:ideal-containments}已给出 \(IJ\subseteq I\cap J\)。若 \(x\in I\cap J\)，取 \(a+b=1\)、\(a\in I,b\in J\)，则 \(x=xa+xb\in IJ\)，这里交换性把 \(xa\) 写成 \(ax\)。

对多个理想，取 \(a_j\in I_j,b_j\in I_r\) 使 \(a_j+b_j=1\)，其中 \(j<r\)。乘积 \(a_1\cdots a_{r-1}\) 属于 \(I_1\cdots I_{r-1}\)，且模 \(I_r\) 同余于 \(1\)。所以 \(I_r+I_1\cdots I_{r-1}=R\)。从 \(r=2\) 起归纳，将两理想结论用于这个乘积理想与 \(I_r\)，即可得到有限情形。
\end{proof}
\term[zhongguo-shengyu-dingli]{中国剩余定理}[Chinese remainder theorem]把这一等式与一个满同态结合起来。

\begin{theorem}[中国剩余定理]\label{b1:thm:crt}
设 \(r\geq1\) 为整数，交换环 \(R\) 的理想 \(I_1,\ldots,I_r\) 两两互素。则映射
\[
\Phi:R\longrightarrow\prod_{i=1}^rR/I_i,\qquad
a\longmapsto(a+I_i)_{i=1}^r
\]
是满环同态，核为 \(\bigcap_iI_i=\prod_iI_i\)，因而诱导环同构
\[
R/\Bigl(\prod_{i=1}^rI_i\Bigr)\cong\prod_{i=1}^rR/I_i.
\]
反过来，若 \(\Phi\) 满射，则这些理想必两两互素。
\end{theorem}
\begin{proof}
逐坐标考察可知 \(\Phi\) 保持运算与单位，其核恰为全部 \(I_i\) 的交。对每个 \(i\ne j\)，取 \(u_{ij}\in I_i,v_{ij}\in I_j\) 使 \(u_{ij}+v_{ij}=1\)，并令
\[
e_i=\prod_{j\ne i}v_{ij}.
\]
空乘积为 \(1\)。于是 \(e_i\equiv1\pmod{I_i}\)，且 \(e_i\in I_j\) 对 \(j\ne i\) 成立。给定任意余数组 \((a_i+I_i)_i\)，元素 \(a=\sum_i a_ie_i\) 的第 \(i\) 个余数恰为 \(a_i+I_i\)，所以 \(\Phi\) 满射。由\cref{b1:lem:comaximal-product}确定核，再对 \(\Phi\) 应用\cref{b1:thm:ring-first-iso}，得到所述同构。

反过来，若 \(\Phi\) 满射，对 \(i\ne j\) 选择一个原像 \(a\)，使 \(a\equiv1\pmod{I_i}\)、\(a\equiv0\pmod{I_j}\)。则 \(1=(1-a)+a\in I_i+I_j\)，故两理想互素。
\end{proof}
有限性使构造中的求和和乘积都在环内有意义。若把定理中的有限族改为无限族，即使理想仍两两互素，满射性也可能失效。例如考虑同余方程组
\[
x\equiv0\pmod2,\qquad
x\equiv1\pmod p\quad\text{对每个奇素数 }p.
\]
其中任意有限多个同余式都由中国剩余定理保证有解。若整数 \(x\) 满足全部同余式，则每个奇素数都整除 \(x-1\)；非零整数只有有限多个素因子，故 \(x-1=0\)。这又与 \(x\equiv0\pmod2\) 矛盾。

\subsection{整数同余方程}
\label{b1:subsec:1004:02}
正整数 \(m,n\) 满足 \(m\mathbb Z+n\mathbb Z=\mathbb Z\)，当且仅当 \(\gcd(m,n)=1\)。若 \(um+vn=1\)，方程
\[
x\equiv a\pmod m,\qquad x\equiv b\pmod n
\]
的解可取 \(x=avn+bum\)。模 \(m\) 第一项同余于 \(a\)、第二项为零；模 \(n\) 恰好相反。\cref{b1:thm:crt}说明全部解组成模 \(mn\) 的唯一剩余类。

\ProseParagraphBreak
例如求 \(x\equiv2\pmod3\)、\(x\equiv3\pmod4\)。由于 \(4-3=1\)，可以取 \(x=2\cdot4+3\cdot(-3)=-1\)，所以答案是 \(x\equiv11\pmod{12}\)。直接核验 \(11\) 的两个余数，可同时检查系数的放置次序。

\begin{proposition}\label{b1:prop:integer-crt-compatible}
设 \(m,n\) 为正整数，\(a,b\in\mathbb Z\)，并令 \(d=\gcd(m,n)\)。同余方程组
\[
x\equiv a\pmod m,\qquad x\equiv b\pmod n
\]
有整数解，当且仅当 \(a\equiv b\pmod d\)。有解时，全部解为模 \(\operatorname{lcm}(m,n)\) 的一个剩余类。
\end{proposition}
\begin{proof}
若有解，\(d\) 同时整除 \(x-a,x-b\)，所以整除 \(a-b\)。反过来，写 \(x=a+mt\)，所求条件成为 \(mt\equiv b-a\pmod n\)。若 \(d\mid b-a\)，约去 \(d\) 后，\(m/d\) 与 \(n/d\) 互素，Bézout 等式使 \(m/d\) 在模 \(n/d\) 下可逆，因而存在 \(t\)。两个解之差必须同时被 \(m,n\) 整除，恰好被它们的最小公倍数整除，且加上这个公倍数保持两个同余式，故得到全部解。
\end{proof}
若 \(n=p_1^{a_1}\cdots p_r^{a_r}\) 为整数的素数幂分解，则
\[
\mathbb Z/n\mathbb Z\cong\prod_{i=1}^r\mathbb Z/p_i^{a_i}\mathbb Z.
\]
这使大模数的乘法、可逆性与幂等元计算分别归结到各个素数幂分量。这里用的是整数的唯一分解；下一章才将元素分解的问题推广到一般整环。

\subsection{幂等元与环的分解}
\label{b1:subsec:1004:03}
\begin{definition}\label{b1:def:idempotent}
设 \(R\) 为环，\(e\in R\)。若 \(e^2=e\)，则称 \(e\) 为\term[midengyuan]{幂等元}[idempotent element]。若一族幂等元 \((e_i)_{i\in I}\) 对任意不同指标 \(i,j\in I\) 都满足 \(e_ie_j=e_je_i=0\)，则称这族幂等元两两正交。
\end{definition}

\begin{proposition}\label{b1:prop:idempotent-product}
设 \(r\geq1\) 为整数。若交换环 \(R\) 有两两正交的幂等元 \(e_1,\ldots,e_r\)，且 \(\sum_i e_i=1\)，则
\[
R\longrightarrow\prod_{i=1}^r e_iR,\qquad
a\longmapsto(e_ia)_i
\]
为环同构。每个 \(e_iR\) 的单位是 \(e_i\)，一般不等于 \(1_R\)。反之，若 \(R_1,\ldots,R_r\) 为交换环，并令
\[
R=\prod_{i=1}^rR_i,
\qquad
e_i=(0,\ldots,0,1_{R_i},0,\ldots,0),
\]
其中 \(1_{R_i}\) 位于第 \(i\) 个坐标，则 \(e_1,\ldots,e_r\) 两两正交、各自幂等，且 \(\sum_i e_i=1_R\)。
\end{proposition}
\begin{proof}
\(e_iR\) 对加法及乘法封闭，\(e_i(e_ia)=e_ia\)，故以 \(e_i\) 为单位。映射保持乘法，因为 \((e_ia)(e_ib)=e_iab\)，并把 \(1\) 送到各分量的单位。若每个 \(e_ia=0\)，则 \(a=\sum_i e_ia=0\)，所以单射。给定 \(x_i\in e_iR\)，令 \(a=\sum_i x_i\)；由正交性，\(e_ia=x_i\)，故满射。逆映射就是各分量求和。有限直积中，第 \(i\) 坐标为单位而其余为零的元素显然满足幂等、正交及和为单位的条件。
\end{proof}
当 \(e_i\ne1_R\) 时，包含映射 \(e_iR\hookrightarrow R\) 不保持单位，因而不是本书约定下的环同态，也不是本书约定下的子环嵌入；\(e_iR\) 在上述分解中以 \(e_i\) 为单位独立成环。

\ProseParagraphBreak
这些幂等元也直接给出中国剩余定理中的理想。对每个 \(i\)，映射
\[
p_i:R\longrightarrow e_iR,\qquad a\longmapsto e_i a
\]
是满环同态，并且 \(p_i(1_R)=e_i=1_{e_iR}\)，所以保持单位。若 \(e_i a=0\)，则 \(a=(1-e_i)a\)；反过来，\(e_i(1-e_i)=0\)。因此
\[
J_i:=\ker p_i=(1-e_i)R,
\qquad e_iR\cong R/J_i.
\]
对于 \(i\ne j\)，由 \(1=(1-e_i)+e_i\) 和 \(e_i=(1-e_j)e_i\in J_j\) 可知 \(J_i+J_j=R\)。若 \(a\in\bigcap_iJ_i\)，则每个 \(e_i a=0\)，于是 \(a=\sum_i e_i a=0\)。故 \(\bigcap_iJ_i=0\)，再由\cref{b1:thm:crt}得到
\[
R\cong\prod_{i=1}^rR/J_i\cong\prod_{i=1}^r e_iR.
\]
这个同构在第 \(i\) 个坐标上的映射就是 \(p_i\)。反过来，从两两互素的理想出发时，\cref{b1:thm:crt}证明中构造的元素 \(e_i\) 未必在原环中幂等，但它们在商环 \(R/\bigcap I_i\) 中的像就是这些坐标幂等元。例如 \(R=\mathbb Z/12\mathbb Z\cong\mathbb Z/3\mathbb Z\times\mathbb Z/4\mathbb Z\) 中，\([4]\) 与 \([9]\) 分别对应 \((1,0)\)、\((0,1)\)：它们的平方各等于自身，乘积为零，和为 \([1]\)。\cref{b1:tab:crt-idempotent-components}列出了相应的分量与投影核。

\begin{table}[htbp]
\centering
\caption{\(\mathbb Z/12\mathbb Z\) 中的幂等元、分量与投影核}
\label{b1:tab:crt-idempotent-components}
\begin{tabularx}{\linewidth}{@{}c>{\centering\arraybackslash}Xcc@{}}
\toprule
幂等元 \(e_i\) & 分量 \(e_iR\)（单位） & 投影核 \(J_i=(1-e_i)R\) & 商环 \(R/J_i\) \\
\midrule
\([4]\) & \(\{[0],[4],[8]\}\)（\([4]\)） & \(([3])\) & \(\mathbb Z/3\mathbb Z\) \\
\([9]\) & \(\{[0],[3],[6],[9]\}\)（\([9]\)） & \(([4])\) & \(\mathbb Z/4\mathbb Z\) \\
\bottomrule
\end{tabularx}
\end{table}

\(e_iR\) 是保留下来的分量，\(J_i=(1-e_i)R\) 则是第 \(i\) 个投影消去的部分；表中两者不是同一个理想，而 \(R/J_i\cong e_iR\)。

\ProseParagraphBreak
整环中 \(e^2=e\) 等价于 \(e(e-1)=0\)，所以只有 \(0,1\) 两个幂等元，不能分成两个非零环的直积。反方向不成立：\(\mathbb Z/4\mathbb Z\) 含零因子，却只有 \([0],[1]\) 两个幂等元。一般非交换环要作上述直积分解，还应要求幂等元与所有元素交换，即为中心幂等元。这样证明中的 \((e_ia)(e_ib)=e_iab\) 才成立。设 \(k\) 为域。在 \(M_2(k)\) 中取 \(e=E_{11}\)、\(f=E_{22}\)，则 \(e^2=e\)、\(f^2=f\)、\(ef=fe=0\) 且 \(e+f=I_2\)，但 \(eE_{12}=E_{12}\ne E_{12}e=0\)。候选坐标映射 \(p(A)=eA\) 也不保持乘法：
\[
p(E_{12}E_{21})=E_{11},\qquad
p(E_{12})p(E_{21})=E_{12}\cdot0=0.
\]
所以正交及和为单位仍不足以取代中心性。


\begin{chaptersummary}
商环要求被抹去的加法子群同时吸收左右乘法，这给出了理想的定义及核的基本性质。交换环的素理想对应整环商，极大理想对应域商；根理想对应没有非零幂零元的商环，准素理想则要求商环中的每个非零零因子都是幂零元。素性与极大性和理想是否非零是不同的问题。理想的和、交、积是三种不同的运算，取根识别哪些元素在商环中成为幂零元，理想商则收集满足指定乘法条件的元素；取零理想作第一项就得到零化子。

\ProseParagraphBreak
中国剩余定理的关键是构造能够分别选择各余数分量的元素；商环的有限直积分解由此转化为正交幂等元的分解。下一章将从主理想的包含关系返回元素整除，考察整数式的分解规律在怎样的环中仍然成立。
\end{chaptersummary}

\BookExercises

\begin{exercise}\label{b1:ex:10-residue-hom}
设 \(m,n\) 是正整数。确定何时存在保单位环同态
\(\mathbb Z/m\mathbb Z\rightarrow\mathbb Z/n\mathbb Z\)，并在存在时给出全部这样的同态、核与像。
\end{exercise}
\begin{solution}
保单位且可加迫使 \([a]_m\mapsto[a]_n\)，所以至多一个同态。由于 \([m]_m=0\)，必须有 \([m]_n=0\)，即 \(n\mid m\)。反过来，若 \(n\mid m\)，同一公式与代表元无关，并保持加法、乘法和单位，故给出同态。它满射，核为
\(\bigl\{\,[kn]_m\mid0\leq k<m/n\,\bigr\}\)，即由 \([n]_m\) 生成的理想，含 \(m/n\) 个元素。\(n=1\) 时目标为零环，公式仍成立；\(m=1\) 时则只能有 \(n=1\)。
\end{solution}

\begin{exercise}\label{b1:ex:10-characteristic-product}
设 \(R=\mathbb Z/4\mathbb Z\times\mathbb Z/6\mathbb Z\)。求 \(R\) 的特征及自然同态 \(\mathbb Z\rightarrow R\) 的核、像，并说明该同态为什么不是满射。
\end{exercise}
\begin{solution}
\(n1_R=0\) 等价于 \(4\mid n\) 且 \(6\mid n\)，最小正解为 \(12\)，故特征为 \(12\)，核为 \(12\mathbb Z\)。\cref{b1:thm:ring-first-iso}给出像同构于 \(\mathbb Z/12\mathbb Z\)，有 \(12\) 个元素；而 \(R\) 有 \(24\) 个元素，故不满射。更具体地，\cref{b1:prop:integer-crt-compatible}说明坐标 \(([a]_4,[b]_6)\) 在像中当且仅当 \(a\equiv b\pmod2\)，因为两模数的最大公因子为二。例如 \(([0]_4,[1]_6)\) 不在像中。环的特征决定这个自然像，却不决定全环的元素个数。
\end{solution}

\begin{exercise}\label{b1:ex:10-matrix-ideals}
设 \(k\) 为域，\(n\geq2\)。证明 \(M_n(k)\) 只有零理想和全环两个双边理想，并说明它既不是整环，也不是除环。
\end{exercise}
\begin{solution}
设非零双边理想 \(I\) 含矩阵 \(A=(a_{ij})\ne0\)，选 \(a_{ij}\ne0\)。对任意 \(r,s\)，矩阵乘法给出
\(E_{ri}AE_{js}=a_{ij}E_{rs}\in I\)。再乘标量矩阵 \(a_{ij}^{-1}I_n\)，得到每个 \(E_{rs}\in I\)。它们的线性组合给出任意矩阵，故 \(I=M_n(k)\)。然而 \(E_{11}E_{22}=0\)，两因子都非零，且 \(E_{11}\) 不可逆。因此在非交换环中，“没有非平凡双边理想”与“每个非零元素都可逆”是两个不同的条件。这里零理想已是极大双边理想，商环仍为 \(M_n(k)\)，却不是除环；\cref{b1:prop:prime-max-quotient}的域商判别确实需要交换性。
\end{solution}

\begin{exercise}\label{b1:ex:10-triangular}
设 \(T\) 为域 \(k\) 上的二阶上三角矩阵环，\(N\) 为严格上三角矩阵的集合。
\begin{enumerate}
\item 求 \(T/N\)、\(T\) 的单位和幂零元，并判断对角矩阵集合是否为 \(T\) 的理想。
\item 令
\[
I=\left\{\begin{pmatrix}a&b\\0&0\end{pmatrix}:a,b\in k\right\},
\qquad
J=\left\{\begin{pmatrix}0&b\\0&c\end{pmatrix}:b,c\in k\right\}.
\]
证明 \(I,J\) 都是 \(T\) 的双边理想，并计算 \(I+J\)、\(I\cap J\)、\(IJ\) 与 \(JI\)。
\end{enumerate}
\end{exercise}
\begin{solution}
1.\ 映射
\(\begin{psmallmatrix}a&b\\0&c\end{psmallmatrix}\mapsto(a,c)\)
是到 \(k\times k\) 的满环同态，核为 \(N\)；\cref{b1:thm:ring-first-iso}给出 \(T/N\cong k\times k\)。当 \(a,c\ne0\) 时，逆矩阵为
\[
\begin{pmatrix}a^{-1}&-a^{-1}bc^{-1}\\0&c^{-1}\end{pmatrix}.
\]
若矩阵可逆，其像 \((a,c)\) 在 \(k\times k\) 中必须可逆，故这也是必要条件。矩阵的 \(r\) 次幂的对角元为 \(a^r,c^r\)，所以幂零必有 \(a=c=0\)；反过来这样的矩阵平方为零。对角矩阵集合含单位但不等于 \(T\)，因而不可能为理想；也可直接计算 \(E_{11}E_{12}=E_{12}\) 看出吸收性失败。

2.\ 取右下对角坐标的同态 \(T\to k\) 的核为 \(I\)，取左上对角坐标的同态之核为 \(J\)，故两者都是双边理想。它们相加含 \(E_{11},E_{12},E_{22}\)，所以 \(I+J=T\)；同时属于两者的矩阵恰为严格上三角矩阵，故 \(I\cap J=N\)。直接相乘可得
\[
\begin{pmatrix}a&b\\0&0\end{pmatrix}
\begin{pmatrix}0&u\\0&c\end{pmatrix}
=\begin{pmatrix}0&au+bc\\0&0\end{pmatrix},
\qquad
\begin{pmatrix}0&u\\0&c\end{pmatrix}
\begin{pmatrix}a&b\\0&0\end{pmatrix}=0.
\]
由于 \(E_{11}E_{12}=E_{12}\)，第一类乘积生成整个 \(N\)，故 \(IJ=N\)，而 \(JI=0\)。这里恰好有 \(IJ=I\cap J\)，但交换两个理想的次序以后，\(J\cap I=N\ne0=JI\)。所以即使两个双边理想之和为全环，它们的乘积仍可能随次序改变，也不能保证乘积等于交；\cref{b1:lem:comaximal-product}中的交换性确实不可省。
\end{solution}

\begin{exercise}\label{b1:ex:10-one-sided-quotient}
设 \(k\) 为域，\(R=M_2(k)\)，\(L\) 是第二列为零的矩阵集合。证明 \(L=RE_{11}\) 是左理想，并用具体代表元说明加法商群 \(R/L\) 不能按 \((A+L)(B+L)=AB+L\) 定义乘法。
\end{exercise}
\begin{solution}
右乘 \(E_{11}\) 保留矩阵第一列并把第二列置零，故 \(L=RE_{11}\)。它对减法封闭，且 \(C(AE_{11})=(CA)E_{11}\in L\)，所以是左理想。\(E_{11}\in L\)，因而 \(E_{11}+L=0+L\)。但是
\(E_{11}E_{12}=E_{12}\notin L\)，而 \(0E_{12}=0\in L\)。同一个左因子陪集的两个代表产生不同的乘积陪集，故乘法不良定义。理想的右吸收性在\cref{b1:prop:quotient-ring}中正是为排除这种现象而要求的。
\end{solution}

\begin{exercise}\label{b1:ex:10-dual-numbers}
设 \(k\) 为域，\(R=k[x]/(x^2)\)，以 \(\varepsilon\) 表示 \(x\) 的像。证明每个元素唯一写成 \(a+b\varepsilon\)，并求 \(R\) 的全部理想、单位和幂零元。利用理想对应，列出 \(k[x]\) 中全部包含 \((x^2)\) 的理想，并判断其中哪些是素理想、哪些是极大理想。
\end{exercise}
\begin{solution}
\((x^2)\) 恰由常数项和一次项均为零的多项式组成，所以每个陪集唯一由 \(a+bx\) 代表，且 \(\varepsilon^2=0\)。若 \(a\ne0\)，则
\[
(a+b\varepsilon)^{-1}=a^{-1}-ba^{-2}\varepsilon.
\]
若 \(a=0\)，元素平方为零，不可能可逆。因此一个真理想中的每个元素均为 \(b\varepsilon\)；只要它含某个 \(b\ne0\) 的元素，乘 \(b^{-1}\) 就含 \(\varepsilon\)，从而等于 \((\varepsilon)\)。全部理想为 \((0),(\varepsilon),R\)。若 \(a+b\varepsilon\) 幂零，其常数项的某次幂为零，故 \(a=0\)；反向已验证。商 \(R/(\varepsilon)\cong k\)，故 \((\varepsilon)\) 是唯一极大理想，且 \((\varepsilon)\ne(0)\)。域的唯一极大理想是 \((0)\)，而这里 \(\varepsilon\ne0\) 且 \(\varepsilon^2=0\)，所以有唯一极大理想并不意味着环是域。

商映射 \(k[x]\to R\) 的核为 \((x^2)\)。由\cref{b1:prop:ring-ideal-correspondence}，包含 \((x^2)\) 的理想恰与上述三个理想对应，依次为 \((x^2),(x),k[x]\)，没有其他情形。其中 \((x)\) 的商环是域 \(k\)，故它既是素理想也是极大理想；\((x^2)\) 不是素理想，因为 \(x^2\in(x^2)\) 而 \(x\notin(x^2)\)。全环不是真理想，故不属于这两类。
\end{solution}

\begin{exercise}\label{b1:ex:10-group-ring-augmentation}
设 \(k\) 为域，\(G\) 为群。在群环 \(k[G]\) 上定义系数和映射
\(\epsilon(\sum_g a_gg)=\sum_g a_g\)。证明它为满环同态，并证明其核作为 \(k[G]\) 的双边理想由全部 \(g-1\)、\(g\in G\) 生成。若 \(G=C_2=\{1,g\}\)，证明 \(k[G]\cong k[t]/(t^2-1)\)。当 \(\operatorname{char}k\ne2\) 时，给出 \(k[G]\cong k\times k\) 的具体同构及其对应的正交幂等元；当 \(\operatorname{char}k=2\) 时，证明增广核的平方为零，并说明群环与 \(k[\varepsilon]/(\varepsilon^2)\) 的关系。
\end{exercise}
\begin{solution}
系数和保持加法与单位；乘积的系数和为
\(\sum_{g,h}a_gb_h=(\sum_ga_g)(\sum_hb_h)\)，所以保持乘法。常数 \(a\cdot1_G\) 映到 \(a\)，故满射。每个 \(g-1\) 属于核，故它们生成的理想包含在核中。反过来，若 \(\sum_g a_g=0\)，则
\[
\sum_g a_gg=\sum_g a_g(g-1),
\]
是这些生成元的有限线性组合，得到反向包含。

现设 \(G=\{1,g\}\)、\(g^2=1\)。代入 \(t\mapsto g\) 给出满环同态 \(k[t]\to k[G]\)，且 \(t^2-1\) 在核中。商环中的关系 \(t^2=1\) 使每个偶次幂化为 \(1\)、每个奇次幂化为 \(t\)，所以每个类都由某个 \(a+bt\) 代表。该表示唯一，因为 \(t^2-1\) 的非零多项式倍数由次数公式具有至少二次，而两个这样的代表之差至多一次。又因 \(1,g\) 是群环的 \(k\)-基，诱导同态 \(k[t]/(t^2-1)\to k[G]\) 也是单射，故为同构。

若 \(\operatorname{char}k\ne2\)，则 \((t-1)+(t+1)=k[t]\)，因为两生成元之差为可逆常数 \(2\)。由\cref{b1:thm:crt}，有具体同构
\[
\begin{aligned}
k[G]&\longrightarrow k\times k,\\
a+bg&\longmapsto(a+b,a-b),
\end{aligned}
\qquad
(u,v)\longmapsto\frac{u+v}{2}+\frac{u-v}{2}g
\]
（右式为其逆）。对应 \((1,0)\)、\((0,1)\) 的正交幂等元分别为 \((1+g)/2\)、\((1-g)/2\)。增广映射是第一个坐标投影，所以其核对应 \(0\times k\)，不是平方为零的理想。

若 \(\operatorname{char}k=2\)，则 \(t+1=t-1\)、\(t^2-1=(t-1)^2\)，两个一次因子不互素，不能作上述直积分解。令 \(\varepsilon=g-1\)，可得 \(k[G]\cong k[\varepsilon]/(\varepsilon^2)\)，其中 \(\varepsilon\ne0\)、\(\varepsilon^2=0\)。增广核为 \(k\varepsilon\)，所以它的平方为零。由\cref{b1:thm:ring-first-iso}，此时的增广商环仍为 \(k\)。这两种情形的分界是域 \(k\) 中的 \(2\) 是否可逆：可逆时两个一次因子互素，商环分裂为两个分量；特征为二时两个因子重合，商环中出现非零幂零元。
\end{solution}

\begin{exercise}\label{b1:ex:10-product-ideals}
设 \(R,S\) 是交换环。证明 \(R\times S\) 的每个理想唯一为 \(I\times J\)，其中 \(I,J\) 分别为两环的理想；进一步确定其全部素理想。
\end{exercise}
\begin{solution}
若 \(K\) 为乘积环的理想，\((a,b)\in K\)，乘 \((1,0)\)、\((0,1)\) 得 \((a,0),(0,b)\in K\)。令
\(I=\bigl\{\,a\mid(a,0)\in K\,\bigr\}\)、\(J=\bigl\{\,b\mid(0,b)\in K\,\bigr\}\)，减法与吸收性直接说明两者为理想，上述拆分及相加给出 \(K=I\times J\)。投影到各坐标又保证唯一性。

若 \(K\) 是素理想，由 \((1,0)(0,1)=0\) 可知至少一个坐标幂等元属于 \(K\)；两者不能同时属于 \(K\)，否则其和 \((1,1)\in K\)。若 \((0,1)\in K\)，则 \(J=S\)，且商环同构于 \(R/I\)，由\cref{b1:prop:prime-max-quotient}知 \(I\) 素；另一种情形给出 \(R\times J\)、\(J\) 素。反过来，这两种形式的商分别是整环，故确为素理想。
\end{solution}

\begin{exercise}\label{b1:ex:10-integer-ideal-operations}
在 \(\mathbb Z\) 中取 \(I=12\mathbb Z,J=18\mathbb Z\)。计算 \(I+J\)、\(I\cap J\)、\(IJ\)、\(I^2\) 及这些理想的根，说明哪些理想不同但根相同。
\end{exercise}
\begin{solution}
因为 \(18-12=6\)，且所有两者的整数线性组合都被六整除，有 \(I+J=6\mathbb Z\)。共同倍数恰为 \(36\) 的倍数，故 \(I\cap J=36\mathbb Z\)。积理想为 \(216\mathbb Z\)，而 \(I^2=144\mathbb Z\)，直接由主理想的乘积定义得到。

所列生成数以及 \(12,18\) 的素因子集合均为 \(\{2,3\}\)。若其中任一个数整除 \(a^r\)，必有 \(2\mid a,3\mid a\)；反过来，若 \(6\mid a\)，取足够大的幂，就能包含该生成数所需的二、三的幂次。因此 \(I,J,I+J,I\cap J,IJ,I^2\) 的根全部为 \(6\mathbb Z\)。这些不同理想保留着不同幂次信息，取根以后该区别消失。
\end{solution}

\begin{exercise}\label{b1:ex:10-integer-colon}
设 \(a,b\) 为正整数，\(d=\gcd(a,b)\)。计算整数环中的理想商 \((a\mathbb Z:b\mathbb Z)\)，并判断何时将它乘以 \(b\mathbb Z\) 能恢复 \(a\mathbb Z\)。再求 \(\mathbb Z/a\mathbb Z\) 中 \([b]_a\) 的零化子及其元素个数。
\end{exercise}
\begin{solution}
写 \(a=da_0,b=db_0\)，其中 \(\gcd(a_0,b_0)=1\)。由\cref{b1:def:ideal-quotient}，\(r\) 属于所求理想当且仅当 \(a\mid rb\)，约去 \(d\) 后等价于 \(a_0\mid rb_0\)。取整数 \(u,v\) 使 \(ua_0+vb_0=1\)，两边乘 \(r\) 可知该条件迫使 \(a_0\mid r\)；反向显然成立。故
\[
(a\mathbb Z:b\mathbb Z)=(a/d)\mathbb Z.
\]
乘回 \(b\mathbb Z\) 后得到 \((ab/d)\mathbb Z\)。由于两个生成数均为正，它等于 \(a\mathbb Z\) 当且仅当 \(b/d=1\)，即 \(b\mid a\)。

在商环中，\([r]_a\) 零化 \([b]_a\) 恰好要求 \(a\mid rb\)，所以零化子由 \([a/d]_a\) 生成，其元素为
\[
\bigl\{\,[j a/d]_a\mid0\leq j<d\,\bigr\}.
\]
这些类互不相同：两者相等要求 \(a\mid(j-j')a/d\)，即 \(d\mid j-j'\)，在给定范围内只能 \(j=j'\)。故零化子恰有 \(d\) 个元素。例如 \(a=24,b=6\) 时，理想商为 \(4\mathbb Z\)，商环中零化子有六个元素。
\end{solution}

\begin{exercise}\label{b1:ex:10-nonprincipal-ideal}
在 \(R=\mathbb Z[x]\) 中令 \(I=(2,x)\)。证明 \(I\) 不是主理想，并计算 \(I^2\)、\(\sqrt{I^2}\) 与 \((I^2:I)\)。
\end{exercise}
\begin{solution}
映射 \(\psi:R\to\mathbb F_2\)、\(f\mapsto[f(0)]_2\) 是满环同态，核为 \(I\)。因此 \(I\) 是真素理想，且 \(\sqrt I=I\)。若 \(I=(f)\)，则由 \(2,x\in I\) 可取 \(g,h\in\mathbb Z[x]\) 使
\[
2=fg,\qquad x=fh.
\]
第一式的两因子均非零，次数公式给出 \(\deg f+\deg g=0\)，所以 \(f,g\) 都是整数常数，且 \(f=\pm1\) 或 \(\pm2\)。再比较第二式的一次项系数，得到 \(1=f h_1\)，其中 \(h_1\in\mathbb Z\) 是 \(h\) 的一次项系数，故只能有 \(f=\pm1\)。于是 \((f)=R\)，与 \(I\ne R\) 矛盾。

由理想积的定义，
\[
I^2=(4,2x,x^2).
\]
其中的多项式恰是常数项被 \(4\) 整除、一次项系数为偶数的多项式。因为 \(I^2\subseteq I=\sqrt I\)，有 \(\sqrt{I^2}\subseteq I\)；反过来，每个 \(f\in I\) 都满足 \(f^2\in I^2\)，故 \(\sqrt{I^2}=I\)。

若 \(f\in(I^2:I)\)，则 \(fx\in I^2\)，由一次项系数的条件得到 \(f(0)\) 为偶数，亦即 \(f\in I\)。反过来，\(I\cdot I=I^2\) 表明每个 \(f\in I\) 都属于 \((I^2:I)\)。所以 \((I^2:I)=I\)。
\end{solution}

\begin{exercise}\label{b1:ex:10-units-lift-nil}
设 \(R\) 为交换环，理想 \(I\) 中每个元素都是幂零元。证明 \(a\in R\) 可逆，当且仅当 \(a+I\) 在 \(R/I\) 中可逆。再给出删去幂零假设后结论失败的例子。
\end{exercise}
\begin{solution}
单位在商同态下仍为单位。反过来，若 \(a+I\) 可逆，选逆的一个代表 \(b\)，则 \(ab=1+n\)，其中 \(n\in I\)。这里只需误差元素 \(n=ab-1\) 自己幂零，不要求存在统一的正整数 \(N\) 使整个理想满足 \(I^N=0\)。对这个 \(n\) 取 \(m\geq1\) 使 \(n^m=0\)，有限几何级数公式给出
\[
(1+n)^{-1}=1-n+n^2-\cdots+(-n)^{m-1}.
\]
于是 \(a\) 的逆为 \(b(1+n)^{-1}\)，左右相乘均为一。若 \(R=\mathbb Z,I=5\mathbb Z,a=2\)，则 \(a+I\) 在商域中可逆，但整数二不可逆；这个理想不满足题设的幂零条件。
\end{solution}

\begin{exercise}\label{b1:ex:10-primary-integers}
对整数 \(n\geq2\)，判定何时 \(n\mathbb Z\) 是准素理想。用 \(n=12\) 写出直接违反准素性条件的两个整数，并用\secref{b1:sec:1003}中的准素商环判别比较 \(\mathbb Z/4\mathbb Z\) 与 \(\mathbb Z/12\mathbb Z\) 中非零零因子的幂零性。
\end{exercise}
\begin{solution}
若 \(n=p^r\) 为素数幂，且 \(ab\in n\mathbb Z\)、\(a\notin n\mathbb Z\)，则 \(p^r\mid ab\) 而 \(p^r\nmid a\)，整数的素数分解迫使 \(p\mid b\)。写 \(b=pc\)，就有 \(b^r=p^rc^r\in n\mathbb Z\)，所以 \(n\mathbb Z\) 准素。

若 \(n\) 不是素数幂，可写 \(n=p^r m\)，其中 \(m>1\)、\(r\geq1\)、\(p\nmid m\)。取 \(a=p^r,b=m\)，则 \(ab=n\in n\mathbb Z\)，但 \(a\notin n\mathbb Z\)，而 \(b\) 的任意正次幂都不被 \(p\) 整除，因而不在 \(n\mathbb Z\) 中。这违反准素性。故 \(n\mathbb Z\) 准素当且仅当 \(n\) 为素数幂。

在 \(n=12\) 时可取 \(a=4,b=3\)：\(4\cdot3\in12\mathbb Z\)，\(4\notin12\mathbb Z\)，且没有 \(3^j\) 是十二的倍数。模四时唯一的非零零因子是 \([2]_4\)，其平方为零；模十二时 \([6]_{12}\) 是平方为零的非零零因子，但 \([3]_{12}\) 被非零的 \([4]_{12}\) 零化而所有正次幂都非零，故零因子并非全都幂零。
\end{solution}

\begin{exercise}\label{b1:ex:10-compatible-congruences}
求全部满足 \(x\equiv3\pmod8\)、\(x\equiv7\pmod{12}\) 的整数，并判断把第二个余数改为 \(5\) 后是否有解。
\end{exercise}
\begin{solution}
两模数的最大公因子为 \(4\)，且 \(3\equiv7\pmod4\)，由\cref{b1:prop:integer-crt-compatible}可解。写 \(x=3+8t\)，代入得 \(8t\equiv4\pmod{12}\)，即 \(2t\equiv1\pmod3\)，所以 \(t\equiv2\pmod3\)。全部解为 \(x=19+24u\)，\(u\in\mathbb Z\)。若第二余数为 \(5\)，则 \(3\not\equiv5\pmod4\)，不满足该命题的必要条件，无解。
\end{solution}

\begin{exercise}\label{b1:ex:10-mod-seventy-two}
求 \(\mathbb Z/72\mathbb Z\) 的全部幂等元与幂零元，并计算其单位个数。
\end{exercise}
\begin{solution}
由\cref{b1:thm:crt}，环同构于 \(\mathbb Z/8\mathbb Z\times\mathbb Z/9\mathbb Z\)，同构为 \([a]_{72}\mapsto([a]_8,[a]_9)\)。若 \(a(a-1)\) 被素数幂 \(p^r\) 整除，因相邻整数互素，该素数幂必全部整除其中一个因子。因此模 \(8\) 和模 \(9\) 各只有 \(0,1\) 两个幂等元；四组坐标提升后恰为
\[
\begin{array}{c|cccc}
\text{模 }72\text{ 的类}&[0]_{72}&[1]_{72}&[9]_{72}&[64]_{72}\\\hline
\text{CRT 坐标}&(0,0)&(1,1)&(1,0)&(0,1)
\end{array}
\]
（各坐标分别在模 \(8\)、模 \(9\) 的环中）。特别地，\([9]_{72}\) 保留模 \(8\) 的分量，\([64]_{72}\) 保留模 \(9\) 的分量。

若 \(72\mid a^m\)，则 \(2,3\) 都整除 \(a\)，故 \(6\mid a\)；反过来 \(6\mid a\) 时 \(72\mid a^3\)。所以幂零元为 \([6j]_{72}\)、\(0\leq j<12\)。单位在模 \(8\) 下有 \(1,3,5,7\) 四个，在模 \(9\) 下有 \(1,2,4,5,7,8\) 六个；直积元素可逆当且仅当两坐标均可逆，因此总数为 \(24\)。
\end{solution}

\begin{exercise}\label{b1:ex:10-comaximal-powers}
设交换环 \(R\) 中 \(I+J=R\)，\(r,s\) 为正整数。证明 \(I^r+J^s=R\)，并给出自然同构
\(R/(I^rJ^s)\cong R/I^r\times R/J^s\)。
\end{exercise}
\begin{solution}
写 \(a+b=1\)，其中 \(a\in I,b\in J\)。展开
\(1=(a+b)^{r+s-1}\)，每一项的 \(a\) 次数至少为 \(r\)，或 \(b\) 次数至少为 \(s\)，所以整个和属于 \(I^r+J^s\)。这证明两幂理想互素。对这两个理想应用\cref{b1:thm:crt}，映射 \(x\mapsto(x+I^r,x+J^s)\) 满射，核为 \(I^r\cap J^s=I^rJ^s\)，故诱导所求同构。此题说明互素关系可以提升到任意正整数幂，而不只是 \(\sqrt{I^r}=\sqrt I\) 所表达的根相同。
\end{solution}

\begin{exercise}\label{b1:ex:10-function-idempotents}
设 \(k\) 为域，\(X\) 为集合。求函数环 \(k^X\) 的全部幂等元。对 \(A\subseteq X\)，利用对应的幂等元解释分解 \(k^X\cong k^A\times k^{X\setminus A}\)，并在 \(X\) 有 \(n\) 个元素时求幂等元个数。
\end{exercise}
\begin{solution}
函数 \(e\) 幂等当且仅当每个 \(x\) 都满足 \(e(x)^2=e(x)\)。域中只有零、一两个幂等元，所以 \(e\) 恰为某个子集 \(A\) 的示性函数：在 \(A\) 上为一，其他点为零。对任意 \(f\in k^X\)，逐点相乘得到
\[
(ef)(x)=\begin{cases}f(x),&x\in A,\\0,&x\notin A,\end{cases}
\qquad
((1-e)f)(x)=\begin{cases}0,&x\in A,\\f(x),&x\notin A.\end{cases}
\]
因此这两个投影分别保留 \(A\) 上和其补集上的函数值。\(e\) 与 \(1-e\) 正交且和为一，\cref{b1:prop:idempotent-product}给出 \(k^X\cong e k^X\times(1-e)k^X\)。两个分量分别通过限制到 \(A\) 与 \(X\setminus A\) 识别为题述函数环；逆映射把两个函数在互不相交的定义域上拼成一个函数。因此有限 \(X\) 的幂等元与其全部子集一一对应，共 \(2^n\) 个。空集时函数环为零环，上述计数仍给出唯一元素。
\end{solution}
